% Preambel mit Einstellungen importieren
% Document type and used packages
\documentclass[open=right, % Sorgt für Umbruch bei Chapter (any erzeugt keine Leerseiten) -> Kapitel darf nur auf der rechten Seite beginnen
    paper=A4,               % DIN-A4-Papier
    a4paper,                % DIN-A4-Papier
    12pt,                   % Schriftgöße
    headings=small,         % Kleine Überschriften
    headsepline=true,       % Trennlinie am Kopf der Seite
    footsepline=false,      % Keine Trennlinie am Fuß der Seite
    bibliography=totoc,     % Literaturverzeichnis in das Inhaltsverzeichnis aufnehmen
    twoside=off,            % Für doppelseitigen Druck auf on stellen, off für einseitig
    DIV=7,                  % Verhältnis der Ränder zum bedruckten Bereich
    chapterprefix=false,    % Kapitel x vor dem Kapitelnamen
    numbers=noendperiod,    % kein Punkt am Ende der Nummerierung
    cleardoublepage=plain]{scrbook}

% Pakete einbinden, die benötigt werden
\usepackage{scrlayer-scrpage}
\usepackage[utf8]{inputenc}       % Dateien in UTF-8 benutzen
\usepackage[T1]{fontenc}          % Zeichenkodierung
\usepackage{graphicx}             % Bilder einbinden
\usepackage[main=ngerman, english]{babel}       % Deutsch und Englisch unterstützen
\usepackage{xcolor}               % Color support
\usepackage{amsmath}              % Matheamtische Formeln
\usepackage{amsfonts}             % Mathematische Zeichensätze
\usepackage{amssymb}              % Mathematische Symbole
\usepackage{float}                % Fließende Objekte (Tabellen, Grafiken etc.)
\usepackage{booktabs}             % Korrekter Tabellensatz
\usepackage[printonlyused, withpage, footnote]{acronym}  % Abkürzungsverzeichnis [nur verwendete Abkürzugen]
\usepackage{makeidx}              % Sachregister
\usepackage{listings}             % Source Code listings
\usepackage{listingsutf8}         % Listings in UTF8
\usepackage[hang,font={sf,footnotesize},labelfont={footnotesize,bf}]{caption} % Beschriftungen
\usepackage[scaled]{helvet}       % Schrift Helvetia laden
\usepackage[absolute]{textpos}	  % Absolute Textpositionen (für Deckblatt)
\usepackage{calc}                 % Berechnung von Positionen
\usepackage{blindtext}            % Blindtexte
\usepackage[left=30mm,right=25mm,top=25mm,bottom=25mm]{geometry} % Ränder ändern
\usepackage{setspace}             % Abstände korrigieren
\usepackage{ifthen}               % Logische Bedingungen mit ifthenelse
\usepackage{scrhack}              % Get rid of tocbasic warnings
\usepackage[pagebackref=false,german]{hyperref}  % Hyperlinks
\usepackage[all]{hypcap}          % Korrekte Verlinkung von Floats
\usepackage[autostyle=true,german=quotes]{csquotes}   % Zitate
\usepackage[backend=biber,
  isbn=true,                     % ISBN nicht anzeigen, gleiches geht mit nahezu allen anderen Feldern
  %sortlocale=de_DE,               % Sortierung der Einträge für Deutsch
  sortlocale=en_US,              % Sortierung der Einträge für Englisch
  autocite=inline,                % regelt Aussehen für \autocite (inline=\parancite)
  hyperref=true,                  % Hyperlinks für Ziate
  %style=ieee                     % Zitate als Zahlen [1]
  %style=alphabetic               % Zitate als Kürzel und Jahr [Ein05]
  %style=authoryear                % Zitate Author und Jahr [Einstein (1905)]
  style=LNI
]{biblatex}                       % Literaturverwaltung mit BibLaTeX
\usepackage{rotating}             % Seiten drehen
\usepackage{harveyballs}          % Harveyballs
\usepackage{tcolorbox}
\usepackage[export]{adjustbox}
\usepackage{subcaption}
\usepackage{color}
\usepackage{colortbl}
\usepackage{wrapfig}
\usepackage{todonotes}
\usepackage{tabularx}
\newcolumntype{b}{>{\hsize=1.2\hsize}X}
\newcolumntype{m}{>{\hsize=.5\hsize}X}
\newcolumntype{s}{>{\hsize=.3\hsize}X}
\usepackage{tikz}


\setlength{\bibitemsep}{1em}     % Abstand zwischen den Literaturangaben
\setlength{\bibhang}{2em}        % Einzug nach jeweils erster Zeile

% Trennung von URLs im Literaturverzeichnis (große Werte [> 10000] verhindern die Trennung)
\defcounter{biburlnumpenalty}{10} % Strafe für Trennung in URL nach Zahl
\defcounter{biburlucpenalty}{500}  % Strafe für Trennung in URL nach Großbuchstaben
\defcounter{biburllcpenalty}{500}  % Strafe für Trennung in URL nach Kleinbuchstaben

% Farben definieren
\definecolor{linkblue}{RGB}{0, 0, 100}
\definecolor{linkblack}{RGB}{0, 0, 0}
\definecolor{comment}{RGB}{63, 127, 95}
\definecolor{darkgreen}{RGB}{14, 144, 102}
\definecolor{darkblue}{RGB}{0,0,168}
\definecolor{darkred}{RGB}{128,0,0}
\definecolor{javadoccomment}{RGB}{0,0,240}
\definecolor{Gray}{RGB}{242,242,242}

% Einstellungen für das Hyperlink-Paket
\hypersetup{
    colorlinks=true,      % Farbige links verwenden
%    allcolors=linkblue,
    linktoc=all,          % Links im Inhaltsverzeichnis
    linkcolor=linkblack,  % Querverweise
    citecolor=linkblack,  % Literaturangaben
	filecolor=linkblack,  % Dateilinks
	urlcolor=linkblack    % URLs
}

% Einstellungen für Quelltexte
\definecolor{backcolour}{rgb}{0.95,0.95,0.92}
\definecolor{codegray}{rgb}{0.5,0.5,0.5}
\lstset{
      xleftmargin=0.5cm,
      xrightmargin=0.5cm,
      frame=lines,
      basicstyle=\small\ttfamily,
      keywordstyle=\color{darkgreen},
      identifierstyle=\color{darkblue},
      commentstyle=\color{comment},
      stringstyle=\color{darkred},
      tabsize=2,
      columns=flexible,
      inputencoding=utf8,
      captionpos=b,
      breakautoindent=true,
	  breakindent=2em,
	  breaklines=true,
	  prebreak=,
	  postbreak=,
      numbers=left,                    
      numbersep=5pt,  
      numberstyle=\tiny\color{codegray},  
      showspaces=false,      % Keine Leerzeichensymbole
      showtabs=false,        % Keine Tabsymbole
      showstringspaces=false,% Leerzeichen in Strings
      morecomment=[s][\color{javadoccomment}]{/**}{*/},
      literate={Ö}{{\"O}}1 {Ä}{{\"A}}1 {Ü}{{\"U}}1 {ß}{{\ss}}2 {ü}{{\"u}}1 {ä}{{\"a}}1 {ö}{{\"o}}1
}


\urlstyle{same}

% Einstellungen für Überschriften
\renewcommand*{\chapterformat}{%
  \Large~\thechapter ~   		% Große Schrift ohne Punkt am Ende
  \vspace{0.3cm}               	% Abstand zum Titel des Kapitels
}

% Abstände für die Überschriften setzen
\renewcommand{\chapterheadstartvskip}{\vspace*{2.6cm}}
\renewcommand{\chapterheadendvskip}{\vspace*{1.5cm}}

\RedeclareSectionCommand[
  beforeskip=-1.8\baselineskip,
  afterskip=0.25\baselineskip]{section}

\RedeclareSectionCommand[
  beforeskip=-1.8\baselineskip,
  afterskip=0.15\baselineskip]{subsection}

\RedeclareSectionCommand[
  beforeskip=-1.8\baselineskip,
  afterskip=0.15\baselineskip]{subsubsection}


% In der Kopfzeile nur die kurze Kapitelbezeichnung (ohne Kapitel davor)
\renewcommand*\chaptermarkformat{\thechapter\autodot\enskip}
\automark[chapter]{chapter}

% Einstellungen für Schriftarten
\setkomafont{pagehead}{\normalfont\sffamily}
\setkomafont{pagenumber}{\normalfont\sffamily}
\setkomafont{paragraph}{\sffamily\bfseries\small}
\setkomafont{subsubsection}{\sffamily\itshape\bfseries\small}
\addtokomafont{footnote}{\footnotesize}
\setkomafont{chapter}{\LARGE\selectfont\bfseries}

% Wichtige Abstände
\setlength{\parskip}{0.2cm}  % 2mm Abstand zwischen zwei Absätzen
\setlength{\parindent}{0mm}  % Absätze nicht einziehen
\clubpenalty = 10000         % Keine "Schusterjungen"
\widowpenalty = 10000        % Keine "Hurenkinder"
\displaywidowpenalty = 10000 % Keine "Hurenkinder"
\renewcommand{\footnotesize}{\fontsize{9}{10}\selectfont} % Größe der Fußnoten
\setlength{\footnotesep}{8pt} % Abstand zwischen den Fußnoten

% Index erzeugen
\makeindex

% Einfacher Font-Wechsel über dieses Makro
\newcommand{\changefont}[3]{
\fontfamily{#1} \fontseries{#2} \fontshape{#3} \selectfont}

% Eigenes Makro für Bilder
\newcommand{\bild}[3]{
\begin{figure}[h]
  \centering
  \includegraphics[width=#2]{#1}
  \caption{#3}
  \label{#1}
\end{figure}}

% Wo liegt Sourcecode?
\newcommand{\srcloc}{src/}

% Wo sind die Bilder?
\graphicspath{{bilder/}}

% Makros für typographisch korrekte Abkürzungen
\newcommand{\zb}[0]{z.\,B.\ }
\newcommand{\dahe}[0]{d.\,h.\ }
\newcommand{\ua}[0]{u.\,a.\ }

% Flags für Veröffentlichung und Sperrvermerk
\newboolean{hsmapublizieren}
\newboolean{hsmasperrvermerk}


% Dokumenteninfos importieren
% In docinfo.tex sind Titel, Autor, Abstract zu definieren
% -------------------------------------------------------
% Daten für die Arbeit
% Wenn hier alles korrekt eingetragen wurde, wird das Titelblatt
% automatisch generiert. D.h. die Datei titelblatt.tex muss nicht mehr
% angepasst werden.

\newcommand{\hsmasprache}{en} % de oder en für Deutsch oder Englisch
% Für korrekt sortierte Literatureinträge, noch preambel.tex anpassen
% und zwar bei \usepackage[main=ngerman, english]{babel},
% \usepackage[pagebackref=false,german]{hyperref}
% und \usepackage[autostyle=true,german=quotes]{csquotes}

% Titel der Arbeit auf Deutsch
% \newcommand{\hsmatitelde}{Bias, Generalisierung und Sim-zu-Real Leistung bei der Instanzsegmentierung: Eine Praktische Untersuchung von Menge und Tiefe}
\newcommand{\hsmatitelde}{Untersuchung des Einflusses von Tiefeninformationen und die Anzahl von Formen und Texturen auf den Bias und die Performance in der Instanz Segmentierung}

% Titel der Arbeit auf Englisch
% \newcommand{\hsmatitelen}{Investigation of the Impact of Material and 3D Model Combinations on Instance Segmentation}
% \newcommand{\hsmatitelen}{Assessing Bias Toward Material or Shape in Instance Segmentation}
% \newcommand{\hsmatitelen}{Investigation of Material and Shape Quantities on Instance Segmentation Accuracy}
%\newcommand{\hsmatitelen}{Investigation of Material and Shape Quantities and Their Bias Impact on Sim-to-Real Instance Segmentation Accuracy}
% \newcommand{\hsmatitelen}{Depth Data and Shape-Texture Biases in Instance Segmentation: Including the Quantity of Varying Shapes and Textures}
% \newcommand{\hsmatitelen}{Generalization and Bias of Shape-Texture and Sim-to-Real Performance in Instance Segmentation: An Analysis of Quantity and Depth}
% Shape-Texture Generalization, Bias and Sim-to-Real in Instance Segmentation: An Analysis of Quantity and Depth
% \newcommand{\hsmatitelen}{Bias, Generalization and Sim-to-Real Performance in Instance Segmentation:\\A Practical Investigation of Quantity and Depth}
% \newcommand{\hsmatitelen}{An Investigation of Depth and Shape-Texture Quantity in Instance Segmentation}
% \newcommand{\hsmatitelen}{Investigating the Influence of Depth Information on the Shape-Texture Bias in Instance Segmentation}
\newcommand{\hsmatitelen}{Investigating the Influence of Depth Information and the Amount of Shapes and Textures on Bias and Performance}
% Investigating the Influence of Depth and Shape-Texture Quantity on Bias and Performance


% \newcommand{\hsmatitelen}{Using 3D-Information Improves Shape Awareness}



% Weitere Informationen zur Arbeit
\newcommand{\hsmaort}{Offenburg}    % Ort
\newcommand{\hsmaautorvname}{Tobia} % Vorname(n)
\newcommand{\hsmaautornname}{Ippolito} % Nachname(n)
\newcommand{\hsmadatum}{06 December} % Datum der Abgabe
\newcommand{\hsmajahr}{2024} % Jahr der Abgabe
\newcommand{\hsmafirma}{Optonic GmbH} % Firma bei der die Arbeit durchgeführt wurde
\newcommand{\hsmabetreuer}{Prof. Dr.-Ing. Janis Keuper, Offenburg University of Applied Sciences} % Betreuer an der Hochschule
\newcommand{\hsmazweitkorrektor}{Moritz Sperling, Optonic GmbH} % Betreuer im Unternehmen oder Zweitkorrektor
\newcommand{\hsmafakultaet}{EMI} % Fakultät
\newcommand{\hsmastudiengang}{AKI} % Studiengangsabkürzung. 
% Diese wird in titelblatt.tex definiert. Bisher AI, EI, MK und INFM. Bitte ergänzen.

% Zustimmung zur Veröffentlichung
\setboolean{hsmapublizieren}{true}   % Einer Veröffentlichung wird zugestimmt
\setboolean{hsmasperrvermerk}{false} % Die Arbeit hat keinen Sperrvermerk

% -------------------------------------------------------
% Abstract

% Kurze (maximal halbseitige) Beschreibung, worum es in der Arbeit geht auf Deutsch
%\newcommand{\hsmaabstractde}{Lorem ipsum dolor sit amet, consetetur sadipscing elitr, sed diam nonumy eirmod tempor invidunt ut labore et dolore magna aliquyam erat, sed diam voluptua. At vero eos et accusam et justo duo dolores et ea rebum. Stet clita kasd gubergren, no sea takimata sanctus est Lorem ipsum dolor sit amet. Lorem ipsum dolor sit amet, consetetur sadipscing elitr, sed diam nonumy eirmod tempor invidunt ut labore et dolore magna aliquyam erat, sed diam voluptua. At vero eos et accusam et justo duo dolores et ea rebum. Stet clita kasd gubergren, no sea takimata sanctus est Lorem ipsum dolor sit amet.}

% Kurze (maximal halbseitige) Beschreibung, worum es in der Arbeit geht auf Englisch

%\newcommand{\hsmaabstracten}{coming soon...}



% Literatur-Datenbank
\addbibresource{literatur.bib}   % BibLaTeX-Datei mit Literaturquellen einbinden

\begin{document}
\frontmatter

% Römische Ziffern für die "Front-Matter"
\setcounter{page}{0}
\changefont{ptm}{m}{n}  % Times New Roman für den Fließtext
\renewcommand{\rmdefault}{ptm}

% Titelblatt
% -------------------------------------------------------
% In dieser Datei sollten eigentlich keine Veränderungen mehr
% notwendig sein.
% -------------------------------------------------------

\thispagestyle{empty}

% Fakultät
% -------------------------------------------------------
\ifthenelse{\equal{\hsmafakultaet}{EI}}%
  {\newcommand{\hsmafakultaetlangde}{Fakultät Elektrotechnik und Informationstechnik}%
   \newcommand{\hsmafakultaetlangen}{Department of Electrical Engineering and Computer Science}}{}
\ifthenelse{\equal{\hsmafakultaet}{EMI}}%
{\newcommand{\hsmafakultaetlangde}{Fakultät Elektrotechnik, Medizintechnik und Informatik}%
	\newcommand{\hsmafakultaetlangen}{Department of Electrical Engineering, Medical Engineering and Computer Science}}{}



\ifthenelse{\equal{\hsmastudiengang}{AI}}%
{\newcommand{\hsmastudienganglangde}{Angewandte Informatik}%
	\newcommand{\hsmastudienganglangen}{Applied Computer Science}%
	\newcommand{\hsmatypde}{BACHELORARBEIT}%
	\newcommand{\hsmatypen}{BACHELOR THESIS}%
	\newcommand{\hsmagrad}{\hsmabachelor}}{}
	
\ifthenelse{\equal{\hsmastudiengang}{AKI}}%
{\newcommand{\hsmastudienganglangde}{Angewandte Künstliche Intelligenz}%
	\newcommand{\hsmastudienganglangen}{Applied Artificial Intelligent}%
	\newcommand{\hsmatypde}{BACHELORARBEIT}%
	\newcommand{\hsmatypen}{BACHELOR THESIS}%
	\newcommand{\hsmagrad}{\hsmabachelor}}{}
	
\ifthenelse{\equal{\hsmastudiengang}{EI}}%
{\newcommand{\hsmastudienganglangde}{Elektrotechnik/Informationstechnik}%
	\newcommand{\hsmastudienganglangen}{Electrical Engineering/Information Technology}%
	\newcommand{\hsmatypde}{BACHELORARBEIT}%
	\newcommand{\hsmatypen}{BACHELOR THESIS}%
	\newcommand{\hsmagrad}{\hsmabachelor}}{}

\ifthenelse{\equal{\hsmastudiengang}{MK}}%
{\newcommand{\hsmastudienganglangde}{Mechatronik}%
	\newcommand{\hsmastudienganglangen}{Mechatronics}%
	\newcommand{\hsmatypde}{BACHELORARBEIT}%
	\newcommand{\hsmatypen}{BACHELOR THESIS}%
	\newcommand{\hsmagrad}{\hsmabachelor}}{}

\ifthenelse{\equal{\hsmastudiengang}{INFM}}%
  {\newcommand{\hsmastudienganglangde}{Informatik Master}%
  \newcommand{\hsmastudienganglangen}{Computer Science Master}%
  \newcommand{\hsmatypde}{MASTERARBEIT}%
  \newcommand{\hsmatypen}{MASTER THESIS}%
  \newcommand{\hsmagrad}{\hsmamaster}}{}

\newcommand{\hsmamaster}{Master of Science (M.Sc.)}

\newcommand{\hsmabachelor}{Bachelor of Science (B.Sc.)}


\newcommand{\hsmakoerperschaftde}{Hochschule für Technik, Wirtschaft und Medien Offenburg}
\newcommand{\hsmakoerperschaften}{Offenburg University}

\newcommand{\hsmaautorbib}{\hsmaautornname, \hsmaautorvname} % Autor Nachname, Vorname
\newcommand{\hsmaautor}{\hsmaautorvname \ \hsmaautornname} % Autor Vorname Nachname

\ifthenelse{\equal{\hsmasprache}{de}}%
  {\newcommand{\hsmatyp}{\hsmatypde}%
   \newcommand{\hsmathesistype}{zur Erlangung des akademischen Grades \hsmagrad}%
   \newcommand{\hsmakoerperschaft}{\hsmakoerperschaftde}%
   \newcommand{\hsmastudiengangname}{Studiengang \hsmastudienganglangde}%
   \newcommand{\hsmastudienganglang}{\hsmastudienganglangde}%
   \newcommand{\hsmatitel}{\hsmatitelde}%
   \newcommand{\hsmatutor}{Betreuer}%
   \newcommand{\hsmafakultaetlang}{\hsmafakultaetlangde}%
   \newcommand{\hsmalistoftables}{Tabellenverzeichnis}%
   \newcommand{\hsmalistoffigures}{Abbildungsverzeichnis}%
   \newcommand{\hsmalistings}{Quellcodeverzeichnis}%
   \newcommand{\hsmaindex}{Index}%
   \newcommand{\hsmaabbreviations}{Abkürzungsverzeichnis}%   
   \selectlanguage{ngerman}}%
  {\newcommand{\hsmatyp}{\hsmatypen}%
   \newcommand{\hsmathesistype}{for the acquisition of the academic degree \hsmagrad}%
   \newcommand{\hsmakoerperschaft}{\hsmakoerperschaften}%
   \newcommand{\hsmastudiengangname}{Course of Studies: \hsmastudienganglang}%
   \newcommand{\hsmastudienganglang}{\hsmastudienganglangen}%
   \newcommand{\hsmatitel}{\hsmatitelen}%
   \newcommand{\hsmatutor}{Tutors}
   \newcommand{\hsmafakultaetlang}{\hsmafakultaetlangen}%
   \newcommand{\hsmalistoftables}{List of Tables}%
   \newcommand{\hsmalistoffigures}{List of Figures}%
   \newcommand{\hsmalistings}{Listings}%
   \newcommand{\hsmaindex}{Index}%
   \newcommand{\hsmaabbreviations}{List of Abbreviations}%
   \selectlanguage{english}}%


% Daten in die Standard-Felder von KOMA-Script eintragen
\titlehead{\hsmatyp\ in\  \hsmastudienganglang}
\subject{}
\title{\hsmatitel}
\author{\hsmaauthor}
\date{\small{\hsmadatum}}

% Daten für das fertige PDF-Dokument
\hypersetup{
  pdftitle={\hsmatitel},  % Titel des Dokuments
  pdfauthor={\hsmaautor},              % Autor
  pdfsubject={\hsmatyp\ in\ \hsmastudienganglang},                % Thema
  pdfkeywords={\hsmatitel}         % Schlüsselworte
}

\newlength{\bindekorrektur}
\newlength{\seitenanfang}
\newlength{\seitenbreite}
  
\setlength{\bindekorrektur}{-46mm}   % Korrektur der horizontalen Position
\setlength{\seitenanfang}{0mm}       % Korrektur der vertikalen Position
\setlength{\seitenbreite}{297mm}

%\noindent \includegraphics[width=7cm, left]{hso.png}\hfill \includegraphics[width=2cm, right]{edeka.png} \\
\captionsetup[figure]{labelformat=empty}
\noindent 
\begin{figure}
	%\includegraphics[width=10cm,center]{hso.jpg}
% Wenn ein Unternehmenslogo mit abgedruckt werden soll,
% kann dies wie folgt integriert werden.	
	\begin{subfigure}[b]{0.5\textwidth}
		\includegraphics[width=7cm,left]{hso.jpg}
	\end{subfigure} 
	\begin{subfigure}[b]{0.5\textwidth}
		\centering
		\includegraphics[width=4cm,right]{optonic.png}
	\end{subfigure} 
	\caption[]{}
\end{figure}
\captionsetup[figure]{labelformat=simple}
% Titel der Arbeit
\begin{textblock*}{128mm}(41mm,\seitenanfang + 62mm) % 4,5cm vom linken Rand und 6,0cm vom oberen Rand
  \centering\Large\sffamily
  \vspace{12mm} % Kleiner zusätzlicher Abstand oben für bessere Optik
  \textbf{\hsmatitel}
\end{textblock*}%

% Name
\begin{textblock*}{\seitenbreite}(\bindekorrektur,\seitenanfang + 108mm)
  \centering\large\sffamily
  \hsmaautor
\end{textblock*}

% Thesis
\begin{textblock*}{\seitenbreite}(\bindekorrektur,\seitenanfang + 130mm)
  \centering\large\sffamily
  \textbf{\hsmatyp}\\
  \begin{small}\hsmathesistype \end{small}\\
  \vspace{6mm}
  \hsmastudiengangname
\end{textblock*}

% Fakultät
\begin{textblock*}{\seitenbreite}(\bindekorrektur,\seitenanfang + 165mm)
  \centering\large\sffamily
  \hsmafakultaetlang\\
  \vspace{2mm}
  \hsmakoerperschaft
\end{textblock*}

% Datum
\begin{textblock*}{\seitenbreite}(\bindekorrektur,\seitenanfang + 190mm)
  \centering\large 
  \textsf{\hsmadatum}
\end{textblock*}

% Firma
\begin{textblock*}{\seitenbreite}(\bindekorrektur,\seitenanfang + 215mm)
  \centering\large 
  % \textsf{Durchgeführt bei der Firma \hsmafirma}
  \textsf{Performed at the company \hsmafirma}
\end{textblock*}

% Betreuer
\begin{textblock*}{\seitenbreite}(\bindekorrektur,\seitenanfang + 240mm)
  \centering\large\sffamily
  \hsmatutor \\
  \vspace{2mm}
  \hsmabetreuer\\
  \vspace{2mm}
  \hsmazweitkorrektor
\end{textblock*}

% Bibliographische Informationen
\null\newpage
\thispagestyle{empty}
  
\newcommand{\hsmabibde}{\begin{small}\textbf{\hsmaautorbib}: \\ \hsmatitelde \ / \hsmaautor. \ -- \\ \hsmatypde, \hsmaort : \hsmakoerperschaftde, \hsmajahr. \pageref{lastpage} Seiten.\end{small}}

\newcommand{\hsmabiben}{\begin{small}\textbf{\hsmaautorbib}: \\ \hsmatitelen \ / \hsmaautor. \ -- \\ \hsmatypen, \hsmaort : \hsmakoerperschaften, \hsmajahr. \pageref{lastpage} pages. \end{small}}

\ifthenelse{\equal{\hsmasprache}{de}}%
  {\hsmabibde \\ \vspace{0.5cm} \\ \hsmabiben}
  {\hsmabiben \\ \vspace{0.5cm} \\ \hsmabibde}


%Vorwort
\clearpage\setcounter{page}{1}
\thispagestyle{empty}
\textsf{\large\textbf{Acknowledgment}}

This work would not have been possible without the guidance of Optonic GmbH, especially Moritz Sperling and my university supervisor, Prof. Dr.-Ing. Janis Keuper.\\
The journey was challenging; I faced many difficulties. Including technological complications and personal health issues. \\
Although I often tried to solve problems independently, seeking help sooner eased some challenges. In the end, I could not be happier to have accepted these challenges and navigated through them with the expertise and dedication of my supervisors.


% Erklärung
\clearpage
\thispagestyle{empty}
% \textsf{\large\textbf{Eidesstattliche Erklärung}}

% Hiermit versichere ich eidesstattlich, dass die vorliegende Master-Thesis von mir selbststän-dig und ohne unerlaubte fremde Hilfe angefertigt worden ist, insbesondere, dass ich alle Stel-len, die wörtlich oder annähernd wörtlich oder dem Gedanken nach aus Veröffentlichungen, unveröffentlichten Unterlagen und Gesprächen entnommen worden sind, als solche an den entsprechenden Stellen innerhalb der Arbeit durch Zitate kenntlich gemacht habe, wobei in den Zitaten jeweils der Umfang der entnommenen Originalzitate kenntlich gemacht wurde. Ich bin mir bewusst, dass eine falsche Versicherung rechtliche Folgen haben wird.

\textsf{\large\textbf{Declaration on oath}}

I hereby declare on oath that this Bachelor's thesis has been prepared by me independently and without unauthorized external assistance, in particular, that I have identified all passages taken verbatim or approximately verbatim or in spirit from publications, unpublished documents, and conversations as such at the appropriate places within the thesis utilizing quotations, whereby the scope of the original quotations taken has been indicated in the quotations. I am aware that a false statement will have legal consequences.

\vspace{0.5cm}

\textsf{\large\textbf{AI-Tool Disclaimer}}

%Nowadays AI-Tools are everywhere and it is important to state what is human-made and what not. In reality, the work of AI and humans is woven together, and therefore, it is even more crucial to name and try to differentiate.\\
% This work was written by a human for humans. 
% The AI-Tool Grammarly \cite{Grammarly} was applied to correct grammar and spelling, for appealing language and best output for the readers. The text itself was not the creation of an AI nor AI assisted.\\
% ChatGPT \cite{ChatGPT} was used for supporting during coding to create and debug tasks efficiently. It still remains the handcraft and thinking of the author.
I hereby affirm that the AI tool Grammarly \cite{Grammarly} was used solely to correct grammar, spelling, and to enhance the language for optimal clarity and appeal to the reader. The text itself, however, was not generated by AI, nor was it AI-assisted in its creation. Furthermore, I testify that ChatGPT \cite{ChatGPT} was employed as a supportive tool during the coding process, assisting with the creation and debugging of tasks in an efficient manner. Nevertheless, the work reflects the independent craftsmanship and reasoning of the author.

\ifthenelse{\boolean{hsmapublizieren} \and \not\boolean{hsmasperrvermerk}}%
{
\vspace{0.5cm}
I agree that my work may be published, i.e. that the work may be stored electronically, converted into other formats, made publicly accessible on the servers of Offenburg University of Applied Sciences and distributed via the Internet.  
}{}%
% Ich bin damit einverstanden, dass meine Arbeit veröffentlicht wird, d.\,h. dass die Arbeit elektronisch gespeichert, in andere Formate konvertiert, auf den Servern der Hochschule Offenburg öffentlich zugänglich gemacht und über das Internet verbreitet werden darf.


\vspace{1cm}
\hsmaort, \hsmadatum \\

\vspace{1.2cm}						                                      
\hsmaautor

\ifthenelse{\boolean{hsmasperrvermerk}}%
{%
\vspace{5cm}
\color{red}\textsf{\large\textbf{Sperrvermerk}}

Die vorliegende Abschlussarbeit beinhaltet vertrauliche Informationen und interne Daten des Unternehmens \hsmafirma.
Sie darf aus diesem Grund nur zu Prüfzwecken verwendet und ohne ausdrückliche Genehmigung durch die \hsmafirma weder Dritten zugänglich gemacht, noch ganz oder in Auszügen veröffentlicht werden. Die Sperrfrist endet 5 Jahre Jahre nach dem Einreichen der Arbeit bei der Hochschule Offenburg. Unbeschadet hiervon bleibt die Weitergabe der Arbeit und Einsicht in die Arbeit an die mit der Prüfung befassten Mitarbeiter der Hochschule und Prüfer möglich, die ihrerseits zur Geheimhaltung verpflichtet sind, sowie die Verwendung der Arbeit in eventuellen prüfungsrechtlichen Rechtsschutzverfahren nach Maßgabe der geltenden verwaltungsprozessualen Regeln.
\color{black}
}{}

\cleardoublepage

% Abstract
% \thispagestyle{empty}
% \textsf{\large\textbf{Zusammenfassung}}
% \subsubsection*{\hsmatitelde}\hsmaabstractde
% \clearpage
\thispagestyle{empty}
\textsf{\large\textbf{Abstract}}
\subsubsection*{\hsmatitelen}\hsmaabstracten



% Inhaltsverzeichnis erzeugen
\cleardoublepage
\pdfbookmark{\contentsname}{Contents}
\tableofcontents

% Korrigiert Nummerierung bei mehrseitigem Inhaltsverzeichnis
\cleardoublepage
\newcounter{frontmatterpage}
\setcounter{frontmatterpage}{\value{page}}

% Arabische Zahlen für den Hauptteil
\mainmatter

% Den Hauptteil mit vergrößertem Zeilenabstand setzen
\onehalfspacing

% ------------------------------------------------------------------
% Hauptteil der Arbeit
% \listoftodos

% Instance Segmentation: A brief look at the current state
% The Need for Accurate Instance Segmentation
% INTRODUCTION
\chapter{Introduction}
\label{chap:kapitel1}



	\section{Objective}    % Importance, Value?
	\label{sec:objective-and-importance}	
		Since introducing instance segmentation in the years 2012-2014 \cite{Yang2012}\cite{Silbermann2012}\cite{Hariharan2014}, it has become one of the most important and complex tasks in computer vision \cite{Sharma2022}. \\
		Detecting and separating every foreground object in an image with pixel-wise accuracy can be found in many modern responsibilities. It ranges from grading prostate cancer \cite{Hassan2022}, to understanding cell division, cellular growth, and morphogenesis \cite{Kar2022}, to wildlife monitoring \cite{Haucke2021}, to segmenting unknown marine objects \cite{Hu2024}, to tooth segmentation in dental medicine \cite{Brahmi2023} and much more. The rising use cases also create a need for more precise segmentations. It can be challenging to achieve this goal because many factors, like the domain, data quality, quantity, \ac{dnn} architecture, and data augmentations, can be adjusted.\\
		For example, it is still being determined if the input images should be provided as RGB with depth or as only depth or RGB. First, it seems straightforward that depth is important and helpful information for segmentation, and it showed promising results in many cases \cite{Danielczuk2019}\cite{Xie2021}. However, recent research shows that RGB-only images as input, for instance segmentation, can achieve even better results \cite{Raj2023}. It does not seem very clear and leads to the question of which factors matter for precise generalization and whether depth data could improve this success even more.\\
		In addition, there are common difficulties in the field of instance segmentation, such as cluttered environments with overlapping objects (also called occlusion), novel objects, and challenging materials like translucent or reflective ones.\\
		More clarity about the shape-texture bias in \ac{cnn}-based approaches must also be clarified. Some research suggests a bias towards texture as rewarding \cite{Qiu2024}. At the same time, others recommend a shape-bias \cite{Geihors2019}, and also a debiased approach can be successful in some scenarios \cite{Li2021}\cite{Co2021}\cite{Chung2023}. In turn, research shows that a bias alone cannot be the reason for a good or bad generalization \cite{Gavrikov2024}. It is essential to notice that some factors of these researches vary, like sometimes the focus is on classification, sometimes on \ac{ood} data, but the current state of research still seems to be confusing and have a gap of research and clearness.\\
		Besides generalization optimization and the influence of bias, there is often the challenging transfer from simulation to the real world. The need for simulations exists due to the data quantity that every \ac{dnn} requires to perform precisely and accurately \cite {Uchida2016}\cite{Alzubaidi2021}\cite{Csurka2023}. Labeling segmentation data is time-consuming and costly since every pixel must be labeled. To solve this challenge, using a virtual environment with automatic labeling is much cheaper and faster. However, there is a gap between the synthetic and real-world data. How to bridge from simulation to the real world still remains a problem of research \cite{Doersch2019}.\\
		Lastly, there is not enough qualitative research about the influence of the quantity of different shapes and textures on instance segmentation.\\
		\\
		\textbf{To conclude}, more research needs to be done on shape-texture bias in instance segmentation, and the existing research needs to be more precise and also in context with generalization. Consequently, it is crucial to continue research on generalization, bias, sim-to-real, depth-data, and shape-texture quantity to improve the performance of instance segmentation.
		
	
	
	\section{Core Focus of the Study}    % Research Gap and 
	\label{sec:core-focus-of-the-study}
		Generalization, shape-texture bias, and sim-to-real transfer are all fundamental and complex areas with many research possibilities. This study will focus on depth data and the quantity of shapes and textures concerning these three areas. The main objective of this study is a practical approach to uncovering new and valuable insights.\\
		Understanding the impact of depth data on shape-texture bias, generalization, and sim-to-real performance of a \ac{dnn} is essential. Collecting depth data in real-world scenarios can be challenging, so knowing how much depth data enhances segmentation outcomes can be highly beneficial.\\
		Additionally, new information about the influence of shape-texture quantity could be relevant and valuable. Generating different shapes and textures with high quality and quantity, potentially in a domain-specific manner, is demanding and time-intensive. Determining the minimum necessary variety of shapes and textures for robust generalization can support the efficiency of \ac{dnn} development for instance-segmentation.
		%A novel examination is the combination of depth data and shape-texture quantity on generalization, shape-texture bias, and sim-to-real transfer, which could yield exciting insights.
		
	
	\clearpage
	\section{Methodology Overview}
	\label{sec:methodology-overview}
		This study presents 18 Mask R-CNN \cite{Kaiming2017} models. Nine models were trained on RGB-only data, and nine were trained on RGB depth data.\\
		Each model is trained on another (here proposed) synthetic dataset. The datasets were created in Unreal Engine 5 \cite{Romero2022} and consist of 20.000 RGB, depth, and mask-images. The shapes and textures used vary between the datasets.\\ 
		With these 18 \ac{dnn}s, three different studies were performed for generalization, shape-texture bias, and sim-to-real transfer. A few sets of data that have special and puzzling textures show if the \ac{dnn} learned to prefer shape or texture. This data includes images where the objects have multiple textures and images where the objects all have the same texture. An in-distribution test dataset with a structured combination of novel and known shapes and textures is proposed to test the performance and generalization of shape and texture.\\
		Finally, all trained \ac{dnn}s are tested with their mean \ac{iou} results on two completely different datasets. One of them is an (also here proposed) real-world dataset is used to see the sim-to-real performance in the bin-picking domain. The other sim-to-real test dataset is the OCID dataset \cite{Suchi2019} which belongs to the common indoor objects domain. 
		
	
	
	\section{Scope and Delimitations}
	\label{sec:scope-and-delimitations}
		Investigating generalization, shape-texture bias, and sim-to-real transfer with a focus on depth data and shape-texture quantity must come with reduced scope and delimitations. Otherwise, it would be an excess of work. \\
		First, only Mask R-CNN \cite{Kaiming2017} is used as \ac{dnn} because many assumptions are made on top of research for \ac{cnn}-based approaches.
		Moreover, all used \ac{dnn}s are trained on the same type of datasets, where only the objects' shape, texture, and position vary.
		Both named delimitations are important in excluding other influences from data and model architecture with certainty. Another \ac{dnn} architecture can differ in bias, accuracy, and learning. Mask R-CNN's wide range of use, distribution, and flexibility made it a good choice as \ac{dnn} for this study.
		For the same reason, every \ac{dnn} used the same hyper-parameters for training, which are described in chapter \ref{chap:kapitel4}. Even if there are more beneficial parameters, these parameters must be consistently equal during the studies to eliminate a possible influence.\\
		This work focuses on instance segmentation for bin-picking. Which is a common task of instance segmentation and important for automation\cite{Raj2023}\cite{Danielczuk2019}\cite{Xie2021}. However, this research should also be valid and vital for other instance segmentation approaches since there are only a few bin-picking-specific influences, like cluttered scenes and rather consistent light settings.\\
		Another delimitation is the selection of \textit{different} shapes and textures. To quantify the difference between two shapes, as well as the difference between two textures, is difficult to measure and remains challenging. There are some methods, like the Structural Similarity Index Measure \cite{Wang2004}, but it is not a standard approach. Therefore, subjective selection was utilized to ensure practical feasibility, as objective methods proved insufficiently reliable for this context.
		\\
		Due to the amount of studies and time limitations, a limited test-data size is given.\\
		This study can find indications and chances, while an underlying truth cannot be provided. Again, this is owed to the complexity and dimension of this research.\\
		Also, due to time limitations, a suboptimal number of epochs and other suboptimal train parameter was chosen. More epochs would lead to even better results, but the extra precision should not influence the experiments; still, it is a limitation.
		Next up, the number of different shapes and textures used is limited to a maximum of 160. The rare availability of high-quality shapes and textures justifies this.\\
		The usage of depth information also comes with delimitations. There are many ways to use depth information, as described in section \ref{sec:state-of-the-art}, but this study can only approach one way to use depth, or else the amount of \ac{dnn} and training would increase significantly. This study uses the straightforward approach of adding a fourth channel with depth information.\\
		Synthetic data often comes with quality issues, and this work's proposed datasets are not different. There are some major downsides which should be considered. The render quality should be higher than it is; however, the cause of this discrepancy remains unclear. The spawning objects can also spawn into each other, sometimes leading to unfavorable data. And the normalization of the depth information not always works properly.
		% Lastly, this study must meet the assumption that it will find limited novel knowledge. This study tries to reinforce or question current state-of-the-art research to find practical results. In a broader context, that could provide valuable insights. A profound look is not possible through the broad research question. Still, there are some novel questions that this study wants to answer.
	
		% only depth could be also interesting -> you said it?
	
	
	\section{Key Definitions}
	\label{sec:key-definitions}
		Here are some important definitions for this work to make a clear point and prevent misunderstandings.\\
		\textbf{Instance Segmentation:} Instance Segmentation is the task of finding and labeling every foreground object in an image pixel-wise. The classification of these objects is optional and plays no role in the segmentation. Defining when an object belongs to the foreground and when to the background is domain-specific and can even differ in a domain depends on the project's requirements. This study defines \textit{foreground objects} as objects that are wanted to be used for bin-picking in an industrial environment.\\ %Most likely every object inside a given bin.\\
		\textbf{Bin Picking:} During bin picking, a robot tries to grasp an object from a bin and uses or transports it to another place. Instance Segmentation is needed in bin picking to find all available objects with their exact position in the bin.\\
		\\
		\textbf{Texture:} Texture is an object's visual appearance and pattern. An object can have the texture of a rock, water, salt, skin, piece of wood, steel beam, and much more. Texture falls into the category of local information. It describes specific areas on a detailed level and does not describe the overall shape.\\
		\textbf{Material:} A material describes the physical properties of an object, like color, metalness, roughness, and the height of the surface. If a picture is taken of an object that has a material, on the picture, the texture of the object is visible, which is defined by its material. The properties of a material can be described as number values, like (0, 255, 0) for the color green or 0.8 for a high metalness, or, for more complexity, as textures.\\
		\textbf{Shape:} A Shape defines the geometry of an object. It is also often referred to as mesh or 3D-model. Simple shapes are spheres and rectangles. Shapes are global information. It defines the boundaries of an object and is only available in the broader view.\\
		\textbf{Depth data:} As depth mentioned data, are images that contain information about the distance from every pixel. The images have only one channel. Each pixel value represents the distance of the corresponding point in reality to the camera.\\
		\textbf{Novel Objects:} Novel objects are objects that were not part of the \ac{dnn}'s training. Novel objects are a typical challenge for bin-picking since there are often unique and innovative products in production and automation pipelines.\\
		\textbf{Domain:} Domain refers to the specific environment or context in which a task or system operates. The domain defines the characteristics of the data; the types of objects, conditions, and setups encountered, such as the industrial environment for bin-picking tasks or the medical environment with different lighting and different objects. The domain influences how models are trained and evaluated, as performance can vary across different domains.
		\textbf{Out of distribution data:} Data which differs essentially from the train data of a \ac{dnn}. \acl{ood} can include novel shapes, textures, backgrounds, resolutions, or other conditions not seen during the training.\\
		% For example a new domain with a completely different background, resolution, quality, and quantity is a \ac{ood} data.
		\textbf{In-distribution data:} Data which is similar to the train data of a \ac{dnn}. The same domain, background and same methodology.\\
		\textbf{Generalization:} The \ac{dnn}'s ability to generalize defines how well it learns the underlying task and patterns of the train data. Shape- or texture-generalization, in instance segmentation specifically, is the ability to learn to segment objects with their shape or texture, even if they are novel to the \ac{dnn}. Meaning that the network learned how to extract shape or texture and detect object boundaries with this information. \\
		\textbf{Bias:} A bias shows a preferred, learned strategy to solve a task. It is most likely caused by the train data itself but also influenced by the \ac{dnn} architecture. A shape- or texture bias is present when the \ac{dnn} prefers the global shape information or the local texture information to make mask predictions.
		\clearpage
		\textbf{Sim-to-real:} Sim-to-real is the transfer from synthetic train data to the real world. Synthetic data is used through its cheap and fast production and the challenging collection of labeled real data. The quality of synthetic data fluctuates but is always much less complex than the natural world and does not reproduce reality precisely enough, thus causing a gap. The goal is that the \ac{dnn} generalize well enough to bridge this sim-to-real gap.\\
		\textbf{Occlusions:} If only a part of an object is visible because another object lies on top, it is referred to as occlusion. Occlusion is a common problem in clutter scenes, appearing often in instance segmentation.
		
	
	
	\section{Structure}
	\label{sec:structure}
		In the \hyperref[chap:kapitel2]{next chapter}, a brief look at the current state of research and related work will be done, which leads to the actual thesis.\\
		Chapter \ref{chap:kapitel3} provides a description of used tools and environments for training and data generation. Furthermore, the implementation of the used \ac{dnn} and data generation will be explained.\\
		Proceeding to chapter \ref{chap:kapitel4} where the concept of experiments and measurement of the results is described.\\
		The 'results' chapter \ref{chap:kapitel5} shows the results of the experiments, followed by a discussion and thoughts about these and what maybe limitates the claims.\\
		Ultimately a summary of the foundings is outlined in chapter \ref{chap:kapitel6}.
		





 % Externe Datei einbinden
\chapter{Research Background and the Hypothesis}
\label{chap:kapitel2}



	\section{State of the Art}
	\label{sec:state-of-the-art}



	\section{Related Work}
	\label{sec:related-work}
	
	
	
	\section{Hypothesis Statement}
	\label{sec:hypothesis-statement}	
	
	
	



			 % Externe Datei einbinden
\chapter{Methodology and Implementation}
\label{chap:kapitel3}
	% To find answers about the previous stated hypotheses in section \ref{sec:hypothesis-statement}, an experimental research design was chosen. An experimental research design suits well to this kind of hypothesis, and other designs would not fit such practical questions. 
	This study proposes a new synthetic dataset for instance segmentation, focusing on bin-picking as domain. The dataset uses the Unreal Engine 5.4.4 \cite{ue5} with hyper-realistic textures and shapes to create the data for the experiments as described in section \ref{sec:data}.\\
	To apply instance segmentation the widely used Mask R-CNN \cite{Kaiming2017} from PyTorch \cite{pytorch} is used \ref{sec:ai-model}.\\
	The coding for data preparation, \ac{dnn} training, and inference was done in Visual Studio Code \cite{vscode}, a flexible and feature-rich \ac{ide}. The data generation was programmed in the blueprint system from Unreal Engine 5 \cite{ue5} in the form of visual programming \cite{Romero2022}.\\
	For managing the Python environments, anaconda \cite{anaconda} got used and greatly impacted the workflow.\\
	The training and inference of all \ac{dnn}s and experiments have been done on three separate remote computers. One Linux-based computer with SSH and X11vnc remote connection from Optonic \cite{optonic} with an NVIDIA RTX 4090 as \ac{gpu}. And two Windows-based computers from Shadow-Tech \cite{shadow}, both with an NVIDIA RTX A4500.\\
	Furthermore this study proposes four in-distribution datasets and one real-world dataset for testing performance and sim-to-real ability, as described in more detail in section \ref{sec:test-data}.

	% \section{Tools and Environment}    % Tools and Environment
	% \label{sec:tools-and-environment}
	
	
	
	\section{AI-Model}
	\label{sec:ai-model}
		Starting with a summary about the used \ac{dnn}, Mask R-CNN \cite{Kaiming2017}. Notice that every implementation of Mask R-CNN can vary in a few points, but the core functionality remains the same.
		Mask R-CNN is the name of a \ac{dnn}, which is historical for instance segmentation. It builds up from Faster R-CNN \cite{Ren2016}, uses the ability to detect objects, and adds the creation of masks for every object. It is a widely adopted architecture and shines with its flexible nature and precise masks.\\
		To understand the fundamentals of this \ac{dnn}, it follows brief explanations about the underlying functionality of the Mask R-CNN followed by a summary to understand the whole process.\\
		\textbf{Residual Networks} (short ResNet) are the first part of Mask R-CNN's architecture. The network is a special type of \ac{cnn}s. \ac{cnn} are key architectures in modern computer vision. They detect spatial features and consist of convolution layers, which convolute images using filters or kernels to detect features like edges, texture, or other patterns, and of pooling layers, which reduce the size of the features by taking only the maximal value from every area and uses activation functions to be able to map non-linear functions \cite{Oshea2015}. The ResNet takes on the feature extraction and is most likely pre-trained. It uses skip connections to ensure the ability to generalize well in deeper layers. Output is a dense version of the origin image, which contains the features from it, called feature-map \cite{He2015}. Thus, the ResNet is the so-called backbone of the Mask R-CNN.\\
		Next, these feature maps get extended through a \textbf{\acl{fpn}} (short \ac{fpn}). The \ac{fpn} is used to get the feature maps on different scales to extract features from larger and smaller objects. The output from multiple ResNet layers is used with \ac{cnn}s to calculate new feature maps. In addition, it creates (also with \ac{cnn}s) a hierarchical structure (like a pyramid) with features in different resolutions and semantic \cite{Lin2017}. The \ac{fpn} also belongs to the backbone of the Mask R-CNN, since it also extracts features.\\
		After extracting the feature maps on different scales, a \textbf{\acl{rpn}} (short \ac{rpn}) uses these feature maps to build proposals that could contain objects called \ac{roi}s. First, the \ac{rpn} creates many anchor boxes for every X pixel. Then, a neural network classifies probabilities for every anchor box and only keeps the ones with high probability. Next up, a regression model refines the remaining anchor boxes. At last, the \ac{roi}s get filtered through \ac{nms}, which first filters the boxes with their scores, then removes boxes with too high \ac{iou}. This step is important because one object can have multiple anchor boxes, which is not wanted \cite{Ren2016}.\\
		One step is left to predict the mask, bounding box, and object class. These predictions require a fixed size of input images (feature maps/\ac{roi}s), but the \ac{roi}s can have different sizes, so a \textbf{\ac{roi} Align} is used. The \ac{roi} Align first combines the anchor boxes with the Feature-Maps (here just called \ac{roi}s). \ac{roi} Align is one method to resize images without much information loss, which is vital for further precise processing. It lays grids over the \ac{roi}s and uses precise float values for the positions, which leads to less information loss. Then a bilinear interpolation is applied to get the new resized \ac{roi}s \cite{Kaiming2017}.
		\clearpage
		One of the heads of Mask R-CNN is the {Object classification and bounding box regression}. A \ac{fc} is used for classification and bounding box regression. The \ac{fc} have a softmax layer to predict the class of the \ac{roi}s (often only background and one other object class). Moreover, a linear regression layer improves the bounding boxes for every \ac{roi} \cite{Ren2016}.\\
		% Since neural networks need the same input size and the \ac{roi}s can differ in width and height, the \ac{roi}s need to get resized (aligned) with as few information reductions as possible.
		Lastly, a \textbf{fully convolutional neural network} (short \ac{fcn}) is used for the prediction of the \ac{roi}'s masks \cite{Kang2014}.\\
		\\
		The procedure of instance segmentation with Mask R-CNN in short:
		\begin{enumerate}
			\item Feature-Map creation with Backbone (ResNet + FPN).
			\item \ac{roi} creation with \ac{rpn} including \ac{nms}.
			\item \ac{roi} align, for equal size conditions.
			\item Creation of the box head, using \ac{fc} for object classification and bounding box refinement.
			\item Creation of mask head, using \ac{fcn}.
		\end{enumerate}
		This architecture can be inspected in the \ref{img:maskrcnn} visualization.\\
		More information can be found here: \cite{Kaiming2017}\cite{Ramesh2021}. 
		
		\begin{figure}[h]
			\centering
			\includegraphics[width=\textwidth]{kapitel3/maskrcnn.png}
			\caption[Visualized Architecture of Mask R-CNN by Tobia Ippolito]{Visualized Architecture of Mask R-CNN}
			\label{img:maskrcnn}
		\end{figure}
		
		This work uses the torchvision implementation from PyTorch \cite{Torchvision}. It is a stable, professional implementation and works with modern graphic cards.\\
		Previously, the famous tensorflow implementation from Matterport \cite{Matterport} was used, but the implementation changed due to issues with newer graphic cards. Section \ref{sec:challenges} covers more of the challenges and learnings.
		
		\clearpage
		There was still much coding to make the \ac{dnn} work as intended. The Python file can be controlled by parameters and is well structured.\\
		One important part is the data loader, which first verifies the data, meaning it checks if there is a solution mask for every image and a depth image (if using depth). Finally, it finds the pixel value used for the background.\\
		During the training process, it loads the RGB image, the solution mask, and the depth image optionally. It also extracts a list of objects and bounding boxes from the solution mask. Data augmentations are significant for the robustness of the \ac{dnn}. These are random transformations and manipulation steps, including flipping, cropping, rotating, blurring, and adding noise to the image. A special data augmentation was applied to add more variation to the background since the background has only two different colors. The background augmentation is listed in \ref{lst:bg-augmentation} since it is not a custom data augmentation. The background augmentation can add noise, a checkerboard pattern, a color gradient, or a color shift. The chosen effect will only be applied to the background using the segmentation masks. \\
		Lastly, the data loader must prepare the input data for the right form for training, resizing, and transforming it into a tensor from PyTorch.
		
		\begin{lstlisting}[language=Python,caption=Random Augmentation of the Background using cv2 in Python, label=lst:bg-augmentation]
class Random_Background_Modification:
		...
		
		def __call__(self, images):
				rgb_img, depth_img, mask_img = images
				rgb_img, depth_img, mask_img = pil_to_cv2([rgb_img, depth_img, 
																																		mask_img])
				
				if random.random() < self.probability:
						mode = random.choice(["noise", "checkerboard",     
											  						 "gradient pattern", "color shift"])
						
						if mode == "noise":
								background_pattern = np.random.randint(0, 256, 
																								(self.height, self.width, 3), 
																								dtype=np.uint8)
						elif mode == "checkerboard":
								checker_size = random.choice([5, 10, 25, 50])
								color1 = random.randint(180, 255)
								color1 = [color1, color1, color1]    # Brighter Color
								color2 = random.randint(0, 130)
								color2 = [color2, color2, color2]    # Darker Color
							
								# Create the checkerboard pattern
								background_pattern = np.zeros((self.height, self.width, 3)
																														, dtype=np.uint8)
								for i in range(0, self.height, checker_size):
										for j in range(0, self.width, checker_size):
												color = color1 if (i // checker_size + 
																							j // checker_size) 
																							% 2 == 0 else color2
												background_pattern[i:i+checker_size, 
																							j:j+checker_size] = color
						elif mode == "gradient pattern":
								background_pattern = np.zeros((self.height, self.width, 3)
																														, dtype=np.uint8)
							
								# Generate a gradient
								if random.random() > 0.5:
										for i in range(self.height):
												color_value = int(255 * (i / self.height))
												background_pattern[i, :] = [color_value, 
																												color_value, 
																												color_value]
								else:
										for i in range(self.width):
												color_value = int(255 * (i / self.width))
												background_pattern[:, i] = [color_value, 
																												color_value, 
																												color_value]
						else:
								B, G, R = cv2.split(rgb_img)
							
								# create shift
								add_B = np.full(B.shape, random.randint(10, 150), 
																												dtype=np.uint8)
								add_G = np.full(G.shape, random.randint(10, 150), 
																												dtype=np.uint8)
								add_R = np.full(R.shape, random.randint(10, 150), 
																												dtype=np.uint8)
								
								# make shift
								shifted_B = cv2.add(B, add_B) if random.random() > 0.5 
																								else cv2.subtract(B, add_B)
								
								shifted_G = cv2.add(G, add_G) if random.random() > 0.5 
																								else cv2.subtract(G, add_G)
								
								shifted_R = cv2.add(R, add_R) if random.random() > 0.5 
																								else cv2.subtract(R, add_R)
								
								# apply shift
								background_pattern = cv2.merge((shifted_B, 
																									shifted_G, 
																									shifted_R))
								
						# apply pattern only on background:
						
						# get pattern in right size
						background_pattern = cv2.resize(background_pattern, 
																	(rgb_img.shape[1], rgb_img.shape[0]))
						
						# Create mask for background and objects
						bg_mask = (mask_img == self.bg_value).astype(np.uint8)
						fg_mask = 1 - bg_mask
						
						# Combine the original image and generated pattern
						background_with_pattern = cv2.bitwise_and(background_pattern, 
																													background_pattern, 
																													mask=bg_mask)
						objects_only = cv2.bitwise_and(rgb_img, rgb_img, mask=fg_mask)
						
						# Overlay the generated pattern and the original objects
						result = cv2.add(background_with_pattern, objects_only)
				else:
						result = rgb_img
				
				# Convert back to cv2
				result, depth_img, mask_img = cv2_to_pil([result, 
																										depth_img, 
																										mask_img])
				return result, depth_img, mask_img
		\end{lstlisting}
		
		\clearpage
		The training function needed experiment tracking, logging, printouts, learn rate scheduling, optimization, model loading, and more. For experiment tracking, mlflow and tensorboard were implemented. Both are helpful tools for tracking training and optimizing the training process. For learn rate scheduling, a simple custom scheduler was chosen with warm-up steps to first make small steps in the correct direction and then be able to increase the learning rate. Choosing a high learning rate at the beginning would lead to swinging up weight over-adjustments. \\
		The model loading and creation process was done without effort through the simple and well-functioning PyTorch implementation. The first layer gets optionally adjusted to add a depth channel to the input image. There is also an adjustment at the loss weights to weigh the segmentation more. The \acl{nms} got adjusted due to the clutter environments in bin picking. Listing \ref{lst:model-loading} shows the loading and adjustment. Notice how simple it is to adjust the architecture. The \ac{dnn} loading function also prints out every layer and a summary of the layers, which is a fascinating insight into the architecture.
		
		\begin{lstlisting}[language=Python,caption=Loading function of Mask R-CNN using torchvision, label=lst:model-loading]
def load_maskrcnn(
				weights_path=None, 
				use_4_channels=False, 
				pretrained=True,
				image_mean=[0.485, 0.456, 0.406, 0.5], 
				image_std=[0.229, 0.224, 0.225, 0.5],    
				min_size=1080, 
				max_size=1920, 
				log_path=None, 
				should_log=False, 
				should_print=True):
		
		backbone = resnet_fpn_backbone(backbone_name='resnet50',
																			weights=ResNet50_Weights.IMAGENET1K_V2
																			) 
		model = MaskRCNN(backbone, num_classes=2)  
		
		if use_4_channels:
				# Change the first Conv2d-Layer for 4 Channels
				in_features = model.backbone.body.conv1.in_channels    
				out_features = model.backbone.body.conv1.out_channels
				kernel_size = model.backbone.body.conv1.kernel_size
				stride = model.backbone.body.conv1.stride
				padding = model.backbone.body.conv1.padding
				
				
				# Create new conv layer with 4 channels
				new_conv1 = torch.nn.Conv2d(4, out_features,
																				kernel_size=kernel_size, 
																				stride=stride, 
																				padding=padding)
				
				# copy the existing weights from the first 3 Channels
				with torch.no_grad():
				new_conv1.weight[:, :3, :, :] = model.backbone.body.conv1.weight  # Copy old 3 Channels
				new_conv1.weight[:, 3:, :, :] = model.backbone.body.conv1.weight
																														[:, :1, :, :]
				
				
				model.backbone.body.conv1 = new_conv1
				
				# Modify the transform to handle 4 channels
				model.transform = GeneralizedRCNNTransform(min_size, max_size, 
																										image_mean, image_std)
		
		# adjust loss weights
		model.rpn.rpn_cls_loss_weight = 1.0
		model.rpn.rpn_bbox_loss_weight = 2.0
		model.roi_heads.mask_loss_weight = 2.0
		model.roi_heads.box_loss_weight = 1.0
		model.roi_heads.classification_loss_weight = 1.0
		
		# adjust non-maximum suppression
		model.roi_heads.nms_thresh = 0.4
		model.roi_heads.box_predictor.nms_thresh = 0.4  
		model.roi_heads.mask_predictor.mask_nms_thresh = 0.4
		model.roi_heads.score_thresh = 0.4
		
		
		# load weights
		if weights_path:
				model.load_state_dict(state_dict=torch.load(weights_path, 
																												weights_only=True)) 
		
		
		
		# printing the architecture
		model_str = "Parameter of Mask R-CNN:"
		model_parts = dict()
		[...]
		
		log(log_path, model_str, should_log=should_log, 
																should_print=should_print)
		
		return model
		\end{lstlisting}
		
		\ac{sgd} with Nesterov-Momentum \cite{Botev2016} was used as learn rate optimizer. \ac{sgd} is known for its good generalization and preciseness but also for its slowness. The Nesterov-Momentum was used to increase convergence speed and stabilize the process.\\
		The train loop itself is a typical PyTorch train loop. \\
		Lastly, the inference is essential to use the \ac{dnn}. The same data loader can be used for the inference due to its flexibility. The inference can process multiple images and create optional visualizations from the created mask and the ground truth (if existing). Exceptional is the feature of visualizing in-between inference steps, as the feature maps from the \ac{fpn} or the \ac{roi} align. For that, hooks were implemented, which call a function when the layer runs. The hook function saves the output of the layer in a dictionary. For the visualization, it was tricky to handle all the different sizes. It is clear to see every described step of the network and makes it sufficiently coherent; thus, listing \ref{lst:inference-insight} shows the described code for the hooking and the visualization, and appendix \ref{appendix:inference-insights} shows samples from one example inference with the network insights.\\
		Additionally, given the ground truth, many different metrics can be calculated during the inference. The inference result is saved as a gray image or a numpy array. The whole process has a print-out for information about the progress.\\
		The coding for the training and inference includes about 66 defined functions, 14 defined classes, one global keyword, 33 imports, 794 called functions and classes, 50 for-loops, 223 if statements, 194 bool operations, and 182 arithmetic operations.
	
		\clearpage
	
	
	\section{Train Data}
	\label{sec:data}
		Nine datasets were created with Unreal Engine 5 \cite{ue5} for the experiments. Each has RGB, mask, and depth images; there are 20.000 images per dataset (60.000 images counting RGB, depth and mask separately). All datasets together need 500-600 GB memory space. The resolution of the images is Full-HD (1920 x 1080). One dataset needed about two days to complete the generation process. Appendix \ref{appendix:traindata-examples} holds examples for every generated dataset.\\
		The visual program in Unreal Engine 5 first makes adjustments to the background, bin box, lighting, and camera. Then, random materials and shapes gets combined and spawned in the bin box. The Materials and Shapes are used from Quixel Megascans \cite{Quixel}, a collection of high-quality assets. Samples of these shapes are available at the appendix \ref{appendix:shapes-for-training} and sample materials/textures are viewed in appendix \ref{appendix:materials-for-training}.\\
		An RGB and depth image is rendered and saved using a structural name convention. The depth image is scaled and converted to the bit-depth of 8; where the ground of the box a high value, near or equal 255, has and the closer objects come, the closer they get to 0 (black).\\
		Next, every object gets a unique color and a new material with that color, and the background and bin-box get a black color. The image is now rendered with a base color renderer, which only takes the base color of every object without any exposure or other effect; the result is an RGB image with 0 as the background, and every object has a unique color. This process will be repeated 20.00 times, and the list of materials and shapes is fixed during this time. The list of materials and shapes is adjusted by switching to the next dataset, and the process starts from new.\\
		\\
		It follows a more in-depth description of visual code for data generation. Only the most important parts can be covered. The code is too massive to show everything in this work here.\\
		The data generator has 24 parameters for customization, as listed in appendix \ref{appendix:custom-params-ue5}.\\
		The initialization nodes only run once after starting the data generation process. First, it checks if the previously listed user parameters are correct. Verification is beneficial to detect early problems and mistakes. Then, every variable receives a reset to ensure everything is ready for the data generation. Figure \ref{img:ue5_init} shows this initial code.
		
		\begin{figure}[h]
			\centering
			\includegraphics[width=\textwidth]{kapitel3/init.jpeg}
			\caption[Visual Code for Initialization in Unreal Engine 5 by Tobia Ippolito]{Visual Code for Initialization in Unreal Engine 5}
			\label{img:ue5_init}
		\end{figure}
		
		According to the initialization, the main loop is called every frame. Many variables and branches are used to determine which steps are already made and which should be done next.\\
		The main loops start with changing the background material and the material of the bin box. After that, the lights get adjusted, which consist of 3 rectangular light boxes. Next, three renderers are created and adjusted to the camera: one for RGB capture, one for depth capture, and one for segmentation mask capture. The next step is noteworthy; here, a random amount of objects with random materials and shapes are scaled, created, and spawned by the bin box. If an object falls out of the box, it automatically gets respawned in the box. A RGB image is taken if all objects stop moving or a given time limit is surpassed. It continues with taking a depth image. Before the segmentation mask image can be made, every object gets a new material with a unique plain color, and the background and bin box are assigned to a black material. The segmentation image can now be taken, and the scene is over. Figure \ref{img:ue5_main} shows a small part of the main loop.
		
		\begin{figure}[h]
			\centering
			\includegraphics[width=\textwidth]{kapitel3/main.jpeg}
			\caption[Visual Code for creating a scene and taking an RGB, depth, and mask image in Unreal Engine 5 by Tobia Ippolito]{Visual Code for creating a scene and taking an RGB, depth, and mask image in Unreal Engine 5}
			\label{img:ue5_main}
		\end{figure}
		
		After creating a scene, every important variable is set to the default value, and the image counter is increased. Then, the program checks whether to change the dataset (changing the material/and shape amount) and reset the image counter or continue with the current dataset. This behavior is presented in figure \ref{img:ue5_dataset_loop}.\\
		When switching to another dataset, there is always a selection of materials and shapes for the whole dataset. The selection can be random or ordered, depending on the chosen parameters. The selection process is viewed in figure \ref{img:ue5_material_shape_choice}.
		
		\begin{figure}[h]
			\centering
			\includegraphics[width=\textwidth]{kapitel3/dataset_loop.jpeg}
			\caption[Visual Code for changing the dataset and prepare the next scene in Unreal Engine 5 by Tobia Ippolito]{Visual Code for changing the dataset and prepare the next scene in Unreal Engine 5}
			\label{img:ue5_dataset_loop}
		\end{figure}
		
		\begin{figure}[h]
			\centering
			\includegraphics[width=\textwidth]{kapitel3/shape_material_choice.jpeg}
			\caption[Visual Code for drawing a subset of a given material and shape list in Unreal Engine 5 by Tobia Ippolito]{Visual Code for drawing a subset of a given material and shape list in Unreal Engine 5}
			\label{img:ue5_material_shape_choice}
		\end{figure}
		
		During the data generation process, Python scripts were also highly required for processing shapes and materials, post-processing the segmentation masks to convert them to gray images, and unzipping a large number of folders.\\
		\\
		The nine in this work proposed datasets called \textbf{3xM} ("triple m") for \textbf{m}odel-\textbf{m}aterial-\textbf{m}ixture which refers to the combinations of different unique amounts of shapes and textures.
		

	\section{Test Data}
	\label{sec:test-data}
		Four in-distribution test datasets with 100 RGB, depth, and mask images were created with the methodology from section \ref{sec:data}.
		The datasets differ in the novelness of their shapes and textures. There is one test dataset with known shapes and known textures, one with unknown shapes and known textures, one with known shapes and unknown textures, and one with unknown shapes and unknown textures. This allows to test not only the performance on in-distribution data but also shows the generalization ability of shapes and textures.\\
		\\
		This study also proposes a real-world dataset in the bin-picking domain. The dataset consists of 70 labeled RGB-D images as shown in appendix \ref{appendix:testdata-examples-simtoreal}. All images were made with the \href{https://www.optonic.com/produkte/ensenso/b57/}{B57-2 camera} from Optonic GmbH \cite{optonic} as shown in figure \ref{img:camera} and consist of 2472 x 2064 pixels.
		\begin{figure}[h]
			\centering
			\includegraphics[width=0.6\textwidth]{kapitel3/camera_B57-4.jpg}
			\caption[RGB-D Camera Ensenso B57-2 from Optonic. More informations on: \url{https://www.optonic.com/produkte/ensenso/b57/}]{RGB-D Camera Ensenso B57-2 from Optonic. More informations on: \url{https://www.optonic.com/produkte/ensenso/b57/}}
			\label{img:camera}
		\end{figure}
		\FloatBarrier
		The setup was a simple black bin with a mat (foam-like) inside and the camera attached to a table, as shown in figure \ref{img:real-data-setup}.
		Twenty-two unique objects were used in an organized way to produce the data scenes: twelve industrial parts and ten everyday consumer goods. In total, 62 objects were used. The scenes are organized in cluttered/uncluttered, solo/mixed (only one type of unique object or different unique object types), and industrial/product/industrial \& product.
		\clearpage
		The objects contain:
		\begin{itemize}
			\item 12x Brass Parts
			\item 9x Complex White Plastic Parts
			\item 8x Round White Plastic Parts
			\item 7x Black Metal Mounting Plates
			\item 5x White T-Form Plastic Parts
			\item 4x Chocolate Candies (different shapes)
			\item 3x Smaller Metal Parts
			\item 3x Small Black Metal Mounting Plates
			\item 2x Metal Parts
			\item 1x USB Stick
			\item 1x Hard Disk
			\item 1x Screwdriver
			\item 1x Pen
			\item 1x Tea
			\item 1x Instant Noodles
			\item 1x Nut Mix
			\item 1x Italian Herbs
			\item 1x Tissues
		\end{itemize}
		Figure \ref{img:real-data-objects} shows all used objects.
		\begin{figure}[h]
			\centering
			\includegraphics[width=0.85\textwidth]{kapitel3/real_data_objects.jpg}
			\caption[All used objects for real-world data collection.]{All used objects for real-world data collection.}
			\label{img:real-data-objects}
		\end{figure}
		
		The connection to the camera happened through the network with the NxView Software from Optonic GmbH. Then, every scene was saved in ZIP file format and later loaded as a file camera in NxView, which behaves similarly to a real camera and was used by the Python script in listing \ref{lst:file-camera}. The listing \ref{lst:file-camera} also shows the used scaling for depth information. The depth scaling starts with decreasing all positive values greater than zero to zero since zero should be the ground of the bin (the camera was calibrated before). Next, the depth data is converted from negative to positive space and missing values are replaced with linear interpolation. If there are values greater than 255, a min-max normalization is applied. A stretching method uses the lowest and highest values to cover the whole range of values. Then, the depth data can be converted to integer values with 8-bit depth, which can be saved and used.
		
		\begin{lstlisting}[language=Python,caption=Capturing images from a file camera with depth image scaling, label=lst:file-camera]
import os
import json
import zipfile

import numpy as np
import cv2
from scipy.interpolate import griddata

from nxlib import NxLib, Camera, NxLibItem
from nxlib.context import NxLib, StereoCamera
from nxlib.command import NxLibCommand
from nxlib.constants import *

with open("./params.json", "r") as file:
		params = json.loads(file.read())

os.makedirs("./data", exist_ok=True)
os.makedirs("./data/rgb", exist_ok=True)
os.makedirs("./data/depth", exist_ok=True)

def interpolate_nan(image):
		x, y = np.indices(image.shape)
		nan_mask = np.isnan(image)
		
		# Coordinates of valid (non-NaN) points
		valid_coords = np.array((x[~nan_mask], y[~nan_mask])).T
		# Values of valid points
		valid_values = image[~nan_mask]
		# Coordinates of NaN points
		nan_coords = np.array((x[nan_mask], y[nan_mask])).T
		
		# Interpolate NaN points
		interpolated_values = griddata(valid_coords, valid_values, nan_coords, method='linear')
		
		# Fill the interpolated values back into the image
		result = image.copy()
		result[nan_mask] = interpolated_values
		
		return result

serial_number = "all_scenes_1"

with NxLib(), StereoCamera(serial_number) as camera:
		camera.get_node()[ITM_PARAMETERS].set_json(json.dumps(params[ITM_PARAMETERS]), True)
		
		first_image = None
		counter = 0
		while True:
				camera.capture()
				camera.rectify()
				camera.compute_disparity_map()
				camera.compute_point_map()
				
				# get unique name
				name = "optonic_bin-picking_dataset_00.png"
				name_counter = 1
				while name in os.listdir("./data/rgb"):
						name = f"optonic_bin-picking_dataset_{name_counter:02}.png"
						name_counter += 1
				
				# calc depth map
				points = camera.get_point_map()
				depth_data = points[:,:,2]
				depth_data = np.minimum(depth_data, 0)
				depth_data = depth_data + abs(np.nanmin(depth_data))
				depth_data = interpolate_nan(depth_data)
				depth_data = depth_data.astype(np.int16)
				
				
				if np.nanmax(depth_data) <= 255:
						depth_data = depth_data.astype(np.uint8)
				else:
						depth_data = ((depth_data - np.nanmax(depth_data)) / 
														(np.nanmax(depth_data) - np.nanmin(depth_data))
														).astype(np.uint8)
						depth_data = (depth_data * 255).astype(np.uint8)
				
				min_val = np.min(depth_data)
				max_val = np.max(depth_data)
				
				# Apply linear scaling to stretch the pixel values
				depth_data = ((depth_data - min_val) / (max_val - min_val)) * 255
				depth_data = np.uint8(depth_data)
				
				# get rgb
				image = camera.get_texture()
				image = cv2.cvtColor(image, cv2.COLOR_BGR2RGB)
				
				if first_image is None:
						first_image = image
				elif np.array_equal(first_image, image):
						break
				
				# save images
				cv2.imwrite(f"./data/depth/{name}", depth_data)
				cv2.imwrite(f"./data/rgb/{name}", image)
				print(f"Successfull extracted image {counter+1}")
				counter += 1
		\end{lstlisting}
		
		The author hand-labeled every image using a semi-automatic annotation approach with CVAT \cite{cvat}. All images were pre-labeled with the SAM model \cite{Kirillov2023} and manually corrected and adjusted. The annotation software runs entirely on the web and does not require any installments. This procedure enabled a quick annotation process without any troubles.
	
		
		\begin{figure}[h]
			\centering
			\includegraphics[width=\textwidth]{kapitel3/real_data_setup.jpg}
			\caption[Setup for collecting real-world bin-picking data.]{Setup for collecting real-world bin-picking data.}
			\label{img:real-data-setup}
		\end{figure}





\chapter{Experiment Setup and Description}
\label{chap:kapitel4}

	All three, here covered, experiments are needed to answer the three hypotheses from \ref{sec:hypothesis-statement}. The relation is not: one experiment for one hypothesis. Multiple experiments are needed to answer one hypothesis.\\
	18 different \acl{dnn} are used in the three experiments. Every \ac{dnn} shares the same training conditions and only differ in the used train data.\\
	9 \ac{dnn}s with only RGB data and 9 with RGB and depth data. The 9 datasets have different object compositions and major difference are the shapes and textures in the images. There are datasets with 10, 80 and 160 shapes and textures. All combinations yield in 9 datasets.\\
	It follows a listing of training conditions which keep the same over all trained networks in this work:
	\begin{itemize}
		\item \textbf{\ac{dnn}:} Mask R-CNN from torchvision
		\item \textbf{Epochs:} 50
		\item \textbf{Data amount:} 20000
		\item \textbf{Width and height:} 1920 x 1080
		\item \textbf{Warm-up iterations:} 2000
		\item \textbf{Learning rate:} 0.003
		\item \textbf{Scheduler:} Simple custom scheduler with warm-up and down regulation
		\item \textbf{Optimizer:} \acl{sgd} with nesterov momentum \cite{Botev2016}
		\item \textbf{Momentum:} 0.9
		\item \textbf{Batch size:} 5
		\item \textbf{Shuffle data:} True
		\item \textbf{Data Augmentations:} Random flip, rotation, crop, brightness contrast, noise, blur, scale and background modification
	\end{itemize}
	
	
	% \section{Experiment Design and Metrics}
	% \label{sec:experiment-design}
	
	\section{Shape vs. Texture Attention Test}
	\label{sec:shape-texutre-attention-test}
		The Shape vs. Texture Attention Test tries to figure out whether the \ac{dnn} prefers to use local information (texture) or global information (shape) for decision-making.\\
		Twenty handcrafted exceptional cases were built to achieve that. The environment of these scenes is equal to the train data, and the awareness of shapes and textures varies. Truly special is that the shapes have multiple textures. For us humans, it is obvious to use the global information (shape) to assign the pixels to the objects, but a \ac{dnn} can decide differently \cite{Geirhos2020}\cite{Mohla2022}\cite{Baker2020}.\\
		% The experiment is still influenced by the fact that \ac{dnn} does not only use one of them to make their decision.
		The 20 tests include five images with known shapes and known textures, five with unknown shapes and known textures, five with known shapes and unknown textures, and five with unknown shapes and unknown textures. This combination of known and unknown shapes and textures can help to get more insights and understanding. For example, it could be that a \ac{dnn} has a bias towards texture but also learned some shape information without enough generalization caused by the texture bias.\\
		It is expected that there is a texture bias by the \ac{dnn} trained on only RGB data due to research \cite{Theodoridis2022}. Thus, the created masks should focus on the texture. \\
		The \ac{dnn} trained on RGB with depth are expected to have a slight bias towards shape since the depth data keeps only shape information. The fact that there is more shape information does not necessarily mean that the \ac{dnn} will favor this kind of information or use this information more than without depth data. Nevertheless, the shape information is more accessible with a fourth depth channel, and the \ac{dnn} probably learned to use this information for better segmentation masks.\\
		The expectation for the shape and texture amount is that more unique shapes and fewer unique textures will lead to higher shape awareness (higher bias towards shape) but less than depth information. The same applies to texture. More unique textures and fewer unique shapes in the train data are expected to increase texture awareness (higher bias towards texture). \\
		No metric will be used in this experiment. There will be a visual investigation of whether the results are orientated towards texture or shape.\\
		\\
		All twenty test images are available in the appendix \ref{appendix:testdata-examples-bias}.
	
	
	
	\section{In-Distribution Performance and Generalization Test}
	\label{sec:in-distribution-performance-generalization}
		Four test datasets with 100 images each were created to test the performance in in-distribution data and generalization. The test is designed to investigate the generalization of global information (shape) and local information (texture). This experiment can also hint in which direction the bias of the \ac{dnn} heads.\\
		One dataset uses shapes and textures used in every \ac{dnn}'s training. This will give insights into the baseline accuracy of the networks. There is no influence by unknown shapes or textures; the composition of objects may differ.\\
		One dataset with unknown shapes and unknown textures will show the performance with novel objects.\\
		There is one dataset with new shapes and known textures for testing the generalization of shapes. A dataset with known shapes and unknown textures tests the texture generalization.\\
		These experiments will show if the \ac{dnn}s just learned the train data or a pattern behind it. Do they learn the shapes in the train data or, more generally, use shape information for segmentation? The same applies to texture.\\
		In addition, this test will show the influence of depth and shape/texture on these investigations.\\
		The general performance is expected to be good but with space for improvements, due to the limited train time. The generalization of novel objects is expected to be good; thus, it still is in-distribution data and is expected to generalize enough for such a task. The shape generalization is expected to improve with depth data and more unique shapes in the train data. The expectation for texture generalization is that a higher amount of textures in the train data leads to a higher generalization, and networks using depth data should also lead to better performance due to the smaller expected texture bias.\\
		This experiment will measure the results using the mean \acl{iou}. A mask can be defined as a two-dimensional complex shape. The \ac{iou} is a standard metric in comparing shapes as stated by \textcite{Rezatofighi2019}, \enquote{IoU, also known as Jaccard index, is the most commonly used metric for comparing the similarity between two arbitrary shapes.}. The \ac{iou} is well comprehensible. The \ac{iou} measure how high the overlay of two shapes is by dividing their area of intersection, where the shapes overlay with each other, with their area of union, the summarized area of both shapes.\\
		The resulting mean \ac{iou}s will be compared.\\
		\\
		Five examples from every in-distribution test dataset are shown in appendix \ref{appendix:testdata-examples-in-distribution}.
	
	
	
	\section{Sim-to-Real Performance Test}
	\label{sec:simt-to-real-performance-test}
		Transferring to real-world data is challenging for \ac{dnn}s only trained on synthetic data. This work will check the real-world accuracy of the Mask R-CNNs presented here. It also wants to provide information about the influence of depth and unique shape/texture amount in train data on sim-to-real ability.\\
		The networks used were only trained on the proposed synthetic datasets and never saw real-world data during their training. Thus, real-world data are \acl{ood}. This work uses the OCID-Dataset proposed by \citeauthor{Suchi2019} as a real-world dataset.\\
		One part of this experiment compares the proposed \ac{dnn} to see the influence of depth and unique shape/texture amount towards sim-to-real ability.\\
		Another part is the comparison with results from the "Learning RGB-D Feature Embeddings for Unseen Object Instance Segmentation" \cite{Xiang2021} paper and the comparison to the results from the in-distribution results to see the general sim-to-real ability also in comparison to the performance to the in-distribution performance.\\
		For measuring the sim-to-real ability the \ac{iou} will be described, as described in section \ref{sec:in-distribution-performance-generalization}.\\
		\\
		The OCID-Dataset is shown with five samples in appendix \ref{appendix:testdata-examples-simtoreal}.
	
	
	
	
	
	
	
	
	



\chapter{Results and Discussion}
\label{chap:kapitel5}

	\section{Results}
	\label{sec:results}



	\section{Interpretation of Results}
	\label{sec:interpretation-of-results}
	
	
	
	\section{Hypotheses Confirmation}
	\label{sec:hypotheses-conformation}
	
	
	
	\section{Discussion} % Limitations and Challenges
	\label{sec:doscussion}
	
	
	









\chapter{Conclusion}
\label{chap:kapitel6}



	\section{Summary of Findings}
	\label{sec:summary-of-findings}
	
	
	
	\section{Implications and Future Work}
	\label{sec:implications-and-future-work}
		% ...
		% Investigating other input types, like only depth images, would be a interesting comparison and could give more insights.
		% ...
		% It also cold be interesting to make a similar investigation towards robustness. Robustness is a import factor of a \ac{dnn}, especially in sensitive fields, like autonomous driving, where new factors like new light settings happens all the time and live could depend on the robustness of a \ac{dnn}. 
	
	
	
	
	
	
	



% ------------------------------------------------------------------

\label{lastpage}

% Neue Seite
\cleardoublepage

% Backmatter mit normalem Zeilenabstand setzen
\singlespacing

% Römische Ziffern für die "Back-Matter", fortlaufend mit "Front-Matter"
\pagenumbering{roman}
%\setcounter{page}{\value{frontmatterpage}}
\setcounter{page}{0}
% Abkürzungsverzeichnis
\addchap{\hsmaabbreviations}
\begin{acronym}[IEEE]
	\acro{cd}[CD]{Continuous Delivery}
	\acro{ci}[CI]{Continuous Integration}
\end{acronym}


% Tabellenverzeichnis erzeugen
\cleardoublepage
\phantomsection
\addcontentsline{toc}{chapter}{\hsmalistoftables}
\listoftables

% Abbildungsverzeichnis erzeugen
\cleardoublepage
\phantomsection
\addcontentsline{toc}{chapter}{\hsmalistoffigures}
\listoffigures

% Listingverzeichnis erzeugen
\cleardoublepage
\phantomsection
\addcontentsline{toc}{chapter}{\hsmalistings}
\lstlistoflistings

% Literaturverzeichnis erzeugen
\begin{flushleft}
\printbibliography
\end{flushleft}

% Index ausgeben. Wenn Sie keinen Index haben, entfernen Sie einfach
% diesen Teil.
%\cleardoublepage
%\phantomsection
%\addcontentsline{toc}{chapter}{\hsmaindex}
%\printindex

% Anhang. Wenn Sie keinen Anhang haben, entfernen Sie einfach
% diesen Teil.
\appendix
\chapter{Inference Insights}
\label{appendix:inference-insights}

	This appendix will show one example inference with outputs from every layer. The visualizations are only samples because there would be too many outputs. At the end of this appendix, there will be the essential code to hook the in-between outputs and how to visualize them.\\
	First presented is the input image.
	
	\begin{figure}[H]
		\centering
		\includegraphics[width=\textwidth]{anhang/insights/Insight_3xM_10000_10_80.png}
		\caption[Insight Inference Input]{Insight Inference Input}
	\end{figure}
	\FloatBarrier
	
	\clearpage
	The first output in Mask R-CNN comes from ResNet. ResNet extracts features from the image; thus, it is called the backbone of Mask R-CNN.
	
	\begin{figure}[H]
		\centering
		\includegraphics[width=\textwidth]{anhang/insights/3xM_10000_10_80_01_resnet_layer1_conv1.jpg}
		\caption[ResNet Layer 1 Convolution 1 Output]{ResNet Layer 1 Convolution 1 Output}
	\end{figure}
	\begin{figure}[H]
		\centering
		\includegraphics[width=\textwidth]{anhang/insights/3xM_10000_10_80_02_resnet_layer1_conv2.jpg}
		\caption[ResNet Layer 1 Convolution 2 Output]{ResNet Layer 1 Convolution 2 Output}
	\end{figure}
	\begin{figure}[H]
		\centering
		\includegraphics[width=\textwidth]{anhang/insights/3xM_10000_10_80_03_resnet_layer1_conv3.jpg}
		\caption[ResNet Layer 1 Convolution 3 Output]{ResNet Layer 1 Convolution 3 Output}
	\end{figure}
	
	
	\begin{figure}[H]
		\centering
		\includegraphics[width=\textwidth]{anhang/insights/3xM_10000_10_80_04_resnet_layer2_conv1.jpg}
		\caption[ResNet Layer 2 Convolution 1 Output]{ResNet Layer 2 Convolution 1 Output}
	\end{figure}
	\begin{figure}[H]
		\centering
		\includegraphics[width=\textwidth]{anhang/insights/3xM_10000_10_80_05_resnet_layer2_conv2.jpg}
		\caption[ResNet Layer 2 Convolution 2 Output]{ResNet Layer 2 Convolution 2 Output}
	\end{figure}
	\begin{figure}[H]
		\centering
		\includegraphics[width=\textwidth]{anhang/insights/3xM_10000_10_80_06_resnet_layer2_conv3.jpg}
		\caption[ResNet Layer 2 Convolution 3 Output]{ResNet Layer 2 Convolution 3 Output}
	\end{figure}
	
	
	\begin{figure}[H]
		\centering
		\includegraphics[width=\textwidth]{anhang/insights/3xM_10000_10_80_07_resnet_layer3_conv1.jpg}
		\caption[ResNet Layer 3 Convolution 1 Output]{ResNet Layer 3 Convolution 1 Output}
	\end{figure}
	\begin{figure}[H]
		\centering
		\includegraphics[width=\textwidth]{anhang/insights/3xM_10000_10_80_08_resnet_layer3_conv2.jpg}
		\caption[ResNet Layer 3 Convolution 2 Output]{ResNet Layer 3 Convolution 2 Output}
	\end{figure}
	\begin{figure}[H]
		\centering
		\includegraphics[width=\textwidth]{anhang/insights/3xM_10000_10_80_09_resnet_layer3_conv3.jpg}
		\caption[ResNet Layer 3 Convolution 3 Output]{ResNet Layer 3 Convolution 3 Output}
	\end{figure}
	
	\begin{figure}[H]
		\centering
		\includegraphics[width=\textwidth]{anhang/insights/3xM_10000_10_80_10_resnet_layer4_conv1.jpg}
		\caption[ResNet Layer 4 Convolution 1 Output]{ResNet Layer 4 Convolution 1 Output}
	\end{figure}
	\begin{figure}[H]
		\centering
		\includegraphics[width=\textwidth]{anhang/insights/3xM_10000_10_80_11_resnet_layer4_conv2.jpg}
		\caption[ResNet Layer 4 Convolution 2 Output]{ResNet Layer 4 Convolution 2 Output}
	\end{figure}
	\begin{figure}[H]
		\centering
		\includegraphics[width=\textwidth]{anhang/insights/3xM_10000_10_80_12_resnet_layer4_conv3.jpg}
		\caption[ResNet Layer 4 Convolution 3 Output]{ResNet Layer 4 Convolution 3 Output}
	\end{figure}
	
	
	\FloatBarrier
	\clearpage
	
	After the ResNet comes the \ac{fpn}. The created feature maps from ResNet will be used on different levels to extract extra output. First, it takes the outputs from ResNet, extracts feature maps from the feature maps, and creates different feature maps on different resolution levels called lateral layers.
	
	\begin{figure}[H]
		\centering
		\includegraphics[width=0.95\textwidth]{anhang/insights/3xM_10000_10_80_19_fpn_lateral_layer4.jpg}
		\caption[\ac{fpn} lateral layer of ResNet Layer 4]{\ac{fpn} lateral layer of ResNet Layer 4}
	\end{figure}
	
	\begin{figure}[H]
		\centering
		\includegraphics[width=0.95\textwidth]{anhang/insights/3xM_10000_10_80_21_fpn_lateral_layer3.jpg}
		\caption[\ac{fpn} lateral layer of ResNet Layer 3]{\ac{fpn} lateral layer of ResNet Layer 3}
	\end{figure}
	
	\begin{figure}[H]
		\centering
		\includegraphics[width=0.95\textwidth]{anhang/insights/3xM_10000_10_80_23_fpn_lateral_layer2.jpg}
		\caption[\ac{fpn} lateral layer of ResNet Layer 2]{\ac{fpn} lateral layer of ResNet Layer 2}
	\end{figure}
	
	\begin{figure}[H]
		\centering
		\includegraphics[width=0.95\textwidth]{anhang/insights/3xM_10000_10_80_25_fpn_lateral_layer1.jpg}
		\caption[\ac{fpn} lateral layer of ResNet Layer 1]{\ac{fpn} lateral layer of ResNet Layer 1}
	\end{figure}
	
	\FloatBarrier
	\clearpage
	The \ac{fpn} uses these feature maps to create the final feature maps on different resolution levels and combines features from different levels.
	
	\begin{figure}[H]
		\centering
		\includegraphics[width=\textwidth]{anhang/insights/3xM_10000_10_80_20_fpn_output_layer4.jpg}
		\caption[\ac{fpn} output layer of ResNet Layer 4]{\ac{fpn} output layer of ResNet Layer 4}
	\end{figure}
	
	\begin{figure}[H]
		\centering
		\includegraphics[width=\textwidth]{anhang/insights/3xM_10000_10_80_22_fpn_output_layer3.jpg}
		\caption[\ac{fpn} output layer of ResNet Layer 3]{\ac{fpn} output layer of ResNet Layer 3}
	\end{figure}
	
	\begin{figure}[H]
		\centering
		\includegraphics[width=\textwidth]{anhang/insights/3xM_10000_10_80_24_fpn_output_layer2.jpg}
		\caption[\ac{fpn} output layer of ResNet Layer 2]{\ac{fpn} output layer of ResNet Layer 2}
	\end{figure}
	
	\begin{figure}[H]
		\centering
		\includegraphics[width=\textwidth]{anhang/insights/3xM_10000_10_80_26_fpn_output_layer1.jpg}
		\caption[\ac{fpn} output layer of ResNet Layer 1]{\ac{fpn} output layer of ResNet Layer 1}
	\end{figure}
	
	
	\FloatBarrier
	
	\clearpage
	Afterward, creating feature maps, the \ac{rpn} localizes object proposals and uses \ac{nms} to remove redundant proposals.\\
	The \ac{rpn} head first uses the feature maps to prepare for object detection and creating proposals, areas within potential objects.\\
	The \ac{rpn} object classification predicts whether a proposal contains an object.\\
	The \ac{rpn} bounding box predictor defines a bounding box for every proposal. The outputs are \ac{roi}s.
	
	\begin{figure}[H]
		\centering
		\includegraphics[width=0.9\textwidth]{anhang/insights/3xM_10000_10_80_27_rpn_head_conv.jpg}
		\caption[\ac{rpn} head convolution. One example Proposal.]{\ac{rpn} head convolution. One example Proposal.}
	\end{figure}
	\begin{figure}[H]
		\centering
		\includegraphics[width=0.9\textwidth]{anhang/insights/3xM_10000_10_80_28_rpn_cls_logits.jpg}
		\caption[\ac{rpn} object classification for one example Proposal]{\ac{rpn} object classification for one example Proposal}
	\end{figure}
	\begin{figure}[H]
		\centering
		\includegraphics[width=0.9\textwidth]{anhang/insights/3xM_10000_10_80_29_rpn_bbox_pred.jpg}
		\caption[\ac{rpn} bounding box prediction for one example Proposal]{\ac{rpn} bounding box prediction for one example Proposal}
	\end{figure}
	
	
	\FloatBarrier
	
	It follows the \ac{roi}-align, where the \ac{roi}s from \ac{rpn} are getting resized.
	
	\begin{figure}[H]
		\centering
		\includegraphics[width=\textwidth]{anhang/insights/3xM_10000_10_80_30_roi_box_pool.jpg}
		\caption[\ac{roi}-align for three example \ac{roi}s]{\ac{roi}-align for three example \ac{roi}s}
	\end{figure}
	
	\FloatBarrier
	\clearpage
	At last, the masks get prepared (refined) by the mask head and then get predicted \ac{roi}-wise.
	
	\begin{figure}[H]
		\centering
		\includegraphics[width=\textwidth]{anhang/insights/3xM_10000_10_80_35_mask_head_conv1.jpg}
		\caption[The mask head uses the standardized \ac{roi}s and refines them for mask prediction]{The mask head uses the standardized \ac{roi}s and refines them for mask prediction}
	\end{figure}
	\begin{figure}[H]
		\centering
		\includegraphics[width=\textwidth]{anhang/insights/3xM_10000_10_80_36_mask_predictor_fcn.jpg}
		\caption[The mask predictor predicts the masks for every polished \ac{roi}]{The mask predictor predicts the masks for every polished \ac{roi}}
	\end{figure}
	
	\FloatBarrier
	
	Hint: The bounding box refinement and object classification are also included at the end, but this insight wanted to focus on mask creation.
	
	
	Following code collect and visualize these fascinating insights.

	\begin{lstlisting}[language=Python,caption=Hooking insight informations from Mask R-CNN, label=lst:inference-insight]
DNN_INSIGHTS = {}

def hook_func(module, input, output, name):
		try:
				DNN_INSIGHTS[name] = {
						'output': output.detach().cpu()
				}
		except AttributeError as e:
				print(f"Error: {e} Data: {input}")

def register_maskrcnn_hooks(model):
		################
		# ResNet hooks #
		################
		model.backbone.body.layer1[0].conv1.register_forward_hook(
				lambda m, i, o: hook_func(m, i, o, 'resnet_layer1_conv1')
		)
		
		model.backbone.body.layer1[0].conv2.register_forward_hook(
				lambda m, i, o: hook_func(m, i, o, 'resnet_layer1_conv2')
		)
		
		model.backbone.body.layer1[0].conv3.register_forward_hook(
				lambda m, i, o: hook_func(m, i, o, 'resnet_layer1_conv3')
		)
		
		model.backbone.body.layer2[0].conv1.register_forward_hook(
				lambda m, i, o: hook_func(m, i, o, 'resnet_layer2_conv1')
		)
		
		model.backbone.body.layer2[0].conv2.register_forward_hook(
				lambda m, i, o: hook_func(m, i, o, 'resnet_layer2_conv2')
		)
		
		model.backbone.body.layer2[0].conv3.register_forward_hook(
				lambda m, i, o: hook_func(m, i, o, 'resnet_layer2_conv3')
		)
		
		model.backbone.body.layer3[0].conv1.register_forward_hook(
				lambda m, i, o: hook_func(m, i, o, 'resnet_layer3_conv1')
		)
		
		model.backbone.body.layer3[0].conv2.register_forward_hook(
				lambda m, i, o: hook_func(m, i, o, 'resnet_layer3_conv2')
		)
		
		model.backbone.body.layer3[0].conv3.register_forward_hook(
				lambda m, i, o: hook_func(m, i, o, 'resnet_layer3_conv3')
		)
		
		model.backbone.body.layer4[0].conv1.register_forward_hook(
				lambda m, i, o: hook_func(m, i, o, 'resnet_layer4_conv1')
		)
		
		model.backbone.body.layer4[0].conv2.register_forward_hook(
				lambda m, i, o: hook_func(m, i, o, 'resnet_layer4_conv2')
		)
		
		model.backbone.body.layer4[0].conv3.register_forward_hook(
				lambda m, i, o: hook_func(m, i, o, 'resnet_layer4_conv3')
		)
		
		#######################################
		# Feature Pyramid Network (FPN) hooks #
		#######################################
		model.backbone.fpn.inner_blocks[0].register_forward_hook(
				lambda m, i, o: hook_func(m, i, o, 'fpn_lateral_layer1')
		)
		
		model.backbone.fpn.inner_blocks[1].register_forward_hook(
				lambda m, i, o: hook_func(m, i, o, 'fpn_lateral_layer2')
		)
		
		model.backbone.fpn.inner_blocks[2].register_forward_hook(
				lambda m, i, o: hook_func(m, i, o, 'fpn_lateral_layer3')
		)
		
		model.backbone.fpn.inner_blocks[3].register_forward_hook(
				lambda m, i, o: hook_func(m, i, o, 'fpn_lateral_layer4')
		)
		
		model.backbone.fpn.layer_blocks[0].register_forward_hook(
				lambda m, i, o: hook_func(m, i, o, 'fpn_output_layer1')
		)
		
		model.backbone.fpn.layer_blocks[1].register_forward_hook(
				lambda m, i, o: hook_func(m, i, o, 'fpn_output_layer2')
		)
		
		model.backbone.fpn.layer_blocks[2].register_forward_hook(
				lambda m, i, o: hook_func(m, i, o, 'fpn_output_layer3')
		)
		
		model.backbone.fpn.layer_blocks[3].register_forward_hook(
				lambda m, i, o: hook_func(m, i, o, 'fpn_output_layer4')
		)
		
		
		#######################################
		# Region Proposal Network (RPN) hooks #
		#######################################
		model.rpn.head.conv.register_forward_hook(
				lambda m, i, o: hook_func(m, i, o, 'rpn_head_conv')
		)
		
		model.rpn.head.cls_logits.register_forward_hook(
				lambda m, i, o: hook_func(m, i, o, 'rpn_cls_logits')
		)
		
		model.rpn.head.bbox_pred.register_forward_hook(
				lambda m, i, o: hook_func(m, i, o, 'rpn_bbox_pred')
		)
		
		###################
		# RoI Align hooks #
		###################
		model.roi_heads.box_roi_pool.register_forward_hook(
				lambda m, i, o: hook_func(m, i, o, 'roi_box_pool')
		)
		
		#############
		# Mask Head #
		#############
		model.roi_heads.mask_head[0].register_forward_hook(
				lambda m, i, o: hook_func(m, i, o, 'mask_head_conv1')
		)
		
		model.roi_heads.mask_predictor.mask_fcn_logits.register_forward_hook(
				lambda m, i, o: hook_func(m, i, o, 'mask_predictor_fcn_logits')
		)




def plot_feature_map(tensor, should_save, save_path, should_show, title="Feature Map"):
		if len(tensor.shape) == 4:
				num_cols = min(tensor.size()[0], 3)
				num_rows = min(tensor.size()[1], 1)
				
				if num_rows > num_cols:
						num_rows = num_cols
				
				fig, all_ax = plt.subplots(num_rows, num_cols, figsize=(15, 10))
				
				
				
				
				
				for cur_col in range(num_cols):
						for cur_row in range(num_rows):
								feature_map_index_1 = cur_col
								feature_map_index_2 = cur_row
								
								if "mask_predictor" in title:
								title = title.replace("_logits", "")
								feature_map_index_1 = feature_map_index_1
								feature_map_index_2 = feature_map_index_2+1
								
								# Extract the feature map for the ith sample
								cur_image = tensor[feature_map_index_1, feature_map_index_2].detach().cpu().numpy()  # Select channel 0, detach from graph
								
								if num_cols > 1 and num_rows == 1:
										ax = all_ax[cur_col]
								elif num_cols == 1 and num_rows == 1:
										ax = all_ax
								elif num_cols == 1 and num_rows > 1:
										ax = all_ax[cur_row]
								elif num_cols > 1 and num_rows > 1:
										ax = all_ax[cur_row][cur_col]
								else:
										raise Exception(f"Error during col: {cur_col}, row: {cur_row}")
								
								ax.imshow(cur_image, cmap='viridis')
								ax.set_title(f'{title} {cur_col+1} {cur_row+1}')
								ax.axis('off')
		elif len(tensor.shape) == 2:
				plot_image = tensor.cpu().numpy()
				plt.figure(figsize=(15, 10))
				plt.imshow(plot_image, cmap='viridis')
				plt.title(title)
				plt.axis('off')
		else:
				raise ValueException(f"Tensor with shape: {tensor.shape} can't be plottet.")
		
		
		if should_save:
				plt.savefig(os.path.join(save_path, f"{title}.jpg"))
		
		if should_show:
				plt.show()
		
		plt.axis('off')
		plt.clf()
		plt.close()


def visualize_insights(insights, should_save, save_path, name, should_show, max_aspect_ratio=5.0, max_cols=3, channel_limit=3, batch_limit=1):
		counter = 1
		for layer_name, data in insights.items():
				try:
						plot_feature_map(data['output'], should_save, save_path, should_show, title=f"{name}_{counter:02}_{layer_name}")
				except Exception as e:
						print(f"Error during insight visualization of: {layer_name} with error: {e} and tensor: {data['output'].size()}")
				counter += 1
	\end{lstlisting}
	



\chapter{Shapes for Training}
\label{appendix:shapes-for-training}

	The amount of unique shapes are one essential part of this work, thus it is important to see some examples of the used shapes.\\
	Here are 20 different shapes, used in the data generation.
	
	\begin{figure}[h]
		\centering
		\begin{subfigure}{0.45\textwidth}
			\centering
			\includegraphics[width=\textwidth]{anhang/shape-1.png}
			\caption[Example of used Shape, created by \cite{Quixel}]{Example of used Shape}
		\end{subfigure}
		\begin{subfigure}{0.45\textwidth}
			\centering
			\includegraphics[width=\textwidth]{anhang/shape-2.png}
			\caption[Example of used Shape, created by \cite{Quixel}]{Example of used Shape}
		\end{subfigure}
		
		\begin{subfigure}{0.45\textwidth}
			\centering
			\includegraphics[width=\textwidth]{anhang/shape-3.png}
			\caption[Example of used Shape, created by \cite{Quixel}]{Example of used Shape}
		\end{subfigure}
		\begin{subfigure}{0.45\textwidth}
			\centering
			\includegraphics[width=\textwidth]{anhang/shape-4.png}
			\caption[Example of used Shape, created by \cite{Quixel}]{Example of used Shape}
		\end{subfigure}
		
		\begin{subfigure}{0.45\textwidth}
			\centering
			\includegraphics[width=\textwidth]{anhang/shape-5.png}
			\caption[Example of used Shape, created by \cite{Quixel}]{Example of used Shape}
		\end{subfigure}
		\begin{subfigure}{0.45\textwidth}
			\centering
			\includegraphics[width=\textwidth]{anhang/shape-6.png}
			\caption[Example of used Shape, created by \cite{Quixel}]{Example of used Shape}
		\end{subfigure}
	
	\end{figure}
	
	
	
	\begin{figure}[h]
		\centering
		\begin{subfigure}{0.45\textwidth}
			\centering
			\includegraphics[width=\textwidth]{anhang/shape-7.png}
			\caption[Example of used Shape, created by \cite{Quixel}]{Example of used Shape}
		\end{subfigure}
		\begin{subfigure}{0.45\textwidth}
			\centering
			\includegraphics[width=\textwidth]{anhang/shape-8.png}
			\caption[Example of used Shape, created by \cite{Quixel}]{Example of used Shape}
		\end{subfigure}
		
		\begin{subfigure}{0.45\textwidth}
			\centering
			\includegraphics[width=\textwidth]{anhang/shape-9.png}
			\caption[Example of used Shape, created by \cite{Quixel}]{Example of used Shape}
		\end{subfigure}
		\begin{subfigure}{0.45\textwidth}
			\centering
			\includegraphics[width=\textwidth]{anhang/shape-10.png}
			\caption[Example of used Shape, created by \cite{Quixel}]{Example of used Shape}
		\end{subfigure}
		
		\begin{subfigure}{0.45\textwidth}
			\centering
			\includegraphics[width=\textwidth]{anhang/shape-11.png}
			\caption[Example of used Shape, created by \cite{Quixel}]{Example of used Shape}
		\end{subfigure}
		\begin{subfigure}{0.45\textwidth}
			\centering
			\includegraphics[width=\textwidth]{anhang/shape-12.png}
			\caption[Example of used Shape, created by \cite{Quixel}]{Example of used Shape}
		\end{subfigure}
		
		\begin{subfigure}{0.45\textwidth}
			\centering
			\includegraphics[width=\textwidth]{anhang/shape-13.png}
			\caption[Example of used Shape, created by \cite{Quixel}]{Example of used Shape}
		\end{subfigure}
		\begin{subfigure}{0.45\textwidth}
			\centering
			\includegraphics[width=\textwidth]{anhang/shape-14.png}
			\caption[Example of used Shape, created by \cite{Quixel}]{Example of used Shape}
		\end{subfigure}
		
	\end{figure}
	
	
	
	\begin{figure}[h]
		\centering
		\begin{subfigure}{0.45\textwidth}
			\centering
			\includegraphics[width=\textwidth]{anhang/shape-15.png}
			\caption[Example of used Shape, created by \cite{Quixel}]{Example of used Shape}
		\end{subfigure}
		\begin{subfigure}{0.45\textwidth}
			\centering
			\includegraphics[width=\textwidth]{anhang/shape-16.png}
			\caption[Example of used Shape, created by \cite{Quixel}]{Example of used Shape}
		\end{subfigure}
	
		\begin{subfigure}{0.45\textwidth}
			\centering
			\includegraphics[width=\textwidth]{anhang/shape-17.png}
			\caption[Example of used Shape, created by \cite{Quixel}]{Example of used Shape}
		\end{subfigure}
		\begin{subfigure}{0.45\textwidth}
			\centering
			\includegraphics[width=\textwidth]{anhang/shape-18.png}
			\caption[Example of used Shape, created by \cite{Quixel}]{Example of used Shape}
		\end{subfigure}
		
		\begin{subfigure}{0.45\textwidth}
			\centering
			\includegraphics[width=\textwidth]{anhang/shape-19.png}
			\caption[Example of used Shape, created by \cite{Quixel}]{Example of used Shape}
		\end{subfigure}
		\begin{subfigure}{0.45\textwidth}
			\centering
			\includegraphics[width=\textwidth]{anhang/shape-20.png}
			\caption[Example of used Shape, created by \cite{Quixel}]{Example of used Shape}
		\end{subfigure}
	
	\end{figure}

\chapter{Materials for Training}
\label{appendix:materials-for-training}

	The amount of unique materials are one essential part of this work, thus it is important to see some examples of the used materials.\\
	Here are 14 different materials, used in the data generation.
	
	\begin{figure}[h]
		\centering
		\begin{subfigure}{0.45\textwidth}
			\centering
			\includegraphics[width=\textwidth]{anhang/material-1.png}
			\caption[Example of used Material, created by \cite{Quixel}]{Example of used Material}
		\end{subfigure}
		\begin{subfigure}{0.45\textwidth}
			\centering
			\includegraphics[width=\textwidth]{anhang/material-2.png}
			\caption[Example of used Material, created by \cite{Quixel}]{Example of used Material}
		\end{subfigure}
		
		\begin{subfigure}{0.45\textwidth}
			\centering
			\includegraphics[width=\textwidth]{anhang/material-3.png}
			\caption[Example of used Material, created by \cite{Quixel}]{Example of used Material}
		\end{subfigure}
		\begin{subfigure}{0.45\textwidth}
			\centering
			\includegraphics[width=\textwidth]{anhang/material-4.png}
			\caption[Example of used Material, created by \cite{Quixel}]{Example of used Material}
		\end{subfigure}
		
		\begin{subfigure}{0.45\textwidth}
			\centering
			\includegraphics[width=\textwidth]{anhang/material-5.png}
			\caption[Example of used Material, created by \cite{Quixel}]{Example of used Material}
		\end{subfigure}
		\begin{subfigure}{0.45\textwidth}
			\centering
			\includegraphics[width=\textwidth]{anhang/material-6.png}
			\caption[Example of used Material, created by \cite{Quixel}]{Example of used Material}
		\end{subfigure}
	\end{figure}
	
	
	
	\begin{figure}[h]
		\centering
		\begin{subfigure}{0.45\textwidth}
			\centering
			\includegraphics[width=\textwidth]{anhang/material-7.png}
			\caption[Example of used Material, created by \cite{Quixel}]{Example of used Material}
		\end{subfigure}
		\begin{subfigure}{0.45\textwidth}
			\centering
			\includegraphics[width=\textwidth]{anhang/material-8.png}
			\caption[Example of used Material, created by \cite{Quixel}]{Example of used Material}
		\end{subfigure}
		
		\begin{subfigure}{0.45\textwidth}
			\centering
			\includegraphics[width=\textwidth]{anhang/material-9.png}
			\caption[Example of used Material, created by \cite{Quixel}]{Example of used Material}
		\end{subfigure}
		\begin{subfigure}{0.45\textwidth}
			\centering
			\includegraphics[width=\textwidth]{anhang/material-10.png}
			\caption[Example of used Material, created by \cite{Quixel}]{Example of used Material}
		\end{subfigure}
		
		\begin{subfigure}{0.45\textwidth}
			\centering
			\includegraphics[width=\textwidth]{anhang/material-11.png}
			\caption[Example of used Material, created by \cite{Quixel}]{Example of used Material}
		\end{subfigure}
		\begin{subfigure}{0.45\textwidth}
			\centering
			\includegraphics[width=\textwidth]{anhang/material-12.png}
			\caption[Example of used Material, created by \cite{Quixel}]{Example of used Material}
		\end{subfigure}
		
		\begin{subfigure}{0.45\textwidth}
			\centering
			\includegraphics[width=\textwidth]{anhang/material-13.png}
			\caption[Example of used Material, created by \cite{Quixel}]{Example of used Material}
		\end{subfigure}
		\begin{subfigure}{0.45\textwidth}
			\centering
			\includegraphics[width=\textwidth]{anhang/material-14.png}
			\caption[Example of used Material, created by \cite{Quixel}]{Example of used Material}
		\end{subfigure}
	\end{figure}
	
	
	\iffalse
	\begin{figure}[h]
		\centering
		\begin{subfigure}{0.45\textwidth}
			\centering
			\includegraphics[width=\textwidth]{anhang/material-15.png}
			\caption[Example of used Material, created by \cite{Quixel}]{Example of used Material}
		\end{subfigure}
		\begin{subfigure}{0.45\textwidth}
			\centering
			\includegraphics[width=\textwidth]{anhang/material-16.png}
			\caption[Example of used Material, created by \cite{Quixel}]{Example of used Material}
		\end{subfigure}
		
		\begin{subfigure}{0.45\textwidth}
			\centering
			\includegraphics[width=\textwidth]{anhang/material-17.png}
			\caption[Example of used Material, created by \cite{Quixel}]{Example of used Material}
		\end{subfigure}
		\begin{subfigure}{0.45\textwidth}
			\centering
			\includegraphics[width=\textwidth]{anhang/material-18.png}
			\caption[Example of used Material, created by \cite{Quixel}]{Example of used Material}
		\end{subfigure}
		
		\begin{subfigure}{0.45\textwidth}
			\centering
			\includegraphics[width=\textwidth]{anhang/material-19.png}
			\caption[Example of used Material, created by \cite{Quixel}]{Example of used Material}
		\end{subfigure}
		\begin{subfigure}{0.45\textwidth}
			\centering
			\includegraphics[width=\textwidth]{anhang/material-20.png}
			\caption[Example of used Material, created by \cite{Quixel}]{Example of used Material}
		\end{subfigure}
	\end{figure}
	\fi

\chapter{Traindata Examples}
\label{appendix:traindata-examples}

	It follows one example per every dataset which this work proposes. Every dataset were generated in Unreal Engine 5 \cite{ue5} and on a NVIDIA RTX A4500.\\
	\\
	One example image for the synthetic dataset with \textbf{10 shapes} and \textbf{10 textures}:
	\begin{figure}[H]
		\centering
		\includegraphics[width=\textwidth]{anhang/example-10-10-1.png}
		\caption[An examle traindata. RGB, depth and the ground truth]{An examle traindata. RGB, depth and the ground truth}
	\end{figure}
	\iffalse
	\begin{figure}[H]
		\centering
		\includegraphics[width=\textwidth]{anhang/example-10-10-2.png}
		\caption[An examle traindata. RGB, depth and the ground truth]{An examle traindata. RGB, depth and the ground truth}
	\end{figure}
	\begin{figure}[H]
		\centering
		\includegraphics[width=\textwidth]{anhang/example-10-10-3.png}
		\caption[An examle traindata. RGB, depth and the ground truth]{An examle traindata. RGB, depth and the ground truth}
	\end{figure}
	\begin{figure}[H]
		\centering
		\includegraphics[width=\textwidth]{anhang/example-10-10-4.png}
		\caption[An examle traindata. RGB, depth and the ground truth]{An examle traindata. RGB, depth and the ground truth}
	\end{figure}
	\begin{figure}[H]
		\centering
		\includegraphics[width=\textwidth]{anhang/example-10-10-5.png}
		\caption[An examle traindata. RGB, depth and the ground truth]{An examle traindata. RGB, depth and the ground truth}
	\end{figure}
	\fi
	
	\FloatBarrier
	% \vspace{3cm}
	% \clearpage
	One example image for the synthetic dataset with \textbf{10 shapes} and \textbf{80 textures}:
	\begin{figure}[H]
		\centering
		\includegraphics[width=\textwidth]{anhang/example-10-80-1.png}
		\caption[An examle traindata. RGB, depth and the ground truth]{An examle traindata. RGB, depth and the ground truth}
	\end{figure}
	\iffalse
	\begin{figure}[H]
		\centering
		\includegraphics[width=\textwidth]{anhang/example-10-80-2.png}
		\caption[An examle traindata. RGB, depth and the ground truth]{An examle traindata. RGB, depth and the ground truth}
	\end{figure}
	\begin{figure}[H]
		\centering
		\includegraphics[width=\textwidth]{anhang/example-10-80-3.png}
		\caption[An examle traindata. RGB, depth and the ground truth]{An examle traindata. RGB, depth and the ground truth}
	\end{figure}
	\begin{figure}[H]
		\centering
		\includegraphics[width=\textwidth]{anhang/example-10-80-4.png}
		\caption[An examle traindata. RGB, depth and the ground truth]{An examle traindata. RGB, depth and the ground truth}
	\end{figure}
	\begin{figure}[H]
		\centering
		\includegraphics[width=\textwidth]{anhang/example-10-80-5.png}
		\caption[An examle traindata. RGB, depth and the ground truth]{An examle traindata. RGB, depth and the ground truth}
	\end{figure}
	\fi
	
	\FloatBarrier
	% \clearpage
	One example image for the synthetic dataset with \textbf{10 shapes} and \textbf{160 textures}:
	\begin{figure}[H]
		\centering
		\includegraphics[width=\textwidth]{anhang/example-10-160-1.png}
		\caption[An examle traindata. RGB, depth and the ground truth]{An examle traindata. RGB, depth and the ground truth}
	\end{figure}
	\iffalse
	\begin{figure}[H]
		\centering
		\includegraphics[width=\textwidth]{anhang/example-10-160-2.png}
		\caption[An examle traindata. RGB, depth and the ground truth]{An examle traindata. RGB, depth and the ground truth}
	\end{figure}
	\begin{figure}[H]
		\centering
		\includegraphics[width=\textwidth]{anhang/example-10-160-3.png}
		\caption[An examle traindata. RGB, depth and the ground truth]{An examle traindata. RGB, depth and the ground truth}
	\end{figure}
	\begin{figure}[H]
		\centering
		\includegraphics[width=\textwidth]{anhang/example-10-160-4.png}
		\caption[An examle traindata. RGB, depth and the ground truth]{An examle traindata. RGB, depth and the ground truth}
	\end{figure}
	\begin{figure}[H]
		\centering
		\includegraphics[width=\textwidth]{anhang/example-10-160-5.png}
		\caption[An examle traindata. RGB, depth and the ground truth]{An examle traindata. RGB, depth and the ground truth}
	\end{figure}
	\fi
	
	\FloatBarrier
	\clearpage
	One example image for the synthetic dataset with \textbf{80 shapes} and \textbf{10 textures}:
	\begin{figure}[H]
		\centering
		\includegraphics[width=\textwidth]{anhang/example-80-10-1.png}
		\caption[An examle traindata. RGB, depth and the ground truth]{An examle traindata. RGB, depth and the ground truth}
	\end{figure}
	\iffalse
	\begin{figure}[H]
		\centering
		\includegraphics[width=\textwidth]{anhang/example-80-10-2.png}
		\caption[An examle traindata. RGB, depth and the ground truth]{An examle traindata. RGB, depth and the ground truth}
	\end{figure}
	\begin{figure}[H]
		\centering
		\includegraphics[width=\textwidth]{anhang/example-80-10-3.png}
		\caption[An examle traindata. RGB, depth and the ground truth]{An examle traindata. RGB, depth and the ground truth}
	\end{figure}
	\begin{figure}[H]
		\centering
		\includegraphics[width=\textwidth]{anhang/example-80-10-4.png}
		\caption[An examle traindata. RGB, depth and the ground truth]{An examle traindata. RGB, depth and the ground truth}
	\end{figure}
	\begin{figure}[H]
		\centering
		\includegraphics[width=\textwidth]{anhang/example-80-10-5.png}
		\caption[An examle traindata. RGB, depth and the ground truth]{An examle traindata. RGB, depth and the ground truth}
	\end{figure}
	\fi
	
	\FloatBarrier
	% \clearpage
	One example image for the synthetic dataset with \textbf{80 shapes} and \textbf{80 textures}:
	\begin{figure}[H]
		\centering
		\includegraphics[width=\textwidth]{anhang/example-80-80-1.png}
		\caption[An examle traindata. RGB, depth and the ground truth]{An examle traindata. RGB, depth and the ground truth}
	\end{figure}
	\iffalse
	\begin{figure}[H]
		\centering
		\includegraphics[width=\textwidth]{anhang/example-80-80-2.png}
		\caption[An examle traindata. RGB, depth and the ground truth]{An examle traindata. RGB, depth and the ground truth}
	\end{figure}
	\begin{figure}[H]
		\centering
		\includegraphics[width=\textwidth]{anhang/example-80-80-3.png}
		\caption[An examle traindata. RGB, depth and the ground truth]{An examle traindata. RGB, depth and the ground truth}
	\end{figure}
	\begin{figure}[H]
		\centering
		\includegraphics[width=\textwidth]{anhang/example-80-80-4.png}
		\caption[An examle traindata. RGB, depth and the ground truth]{An examle traindata. RGB, depth and the ground truth}
	\end{figure}
	\begin{figure}[H]
		\centering
		\includegraphics[width=\textwidth]{anhang/example-80-80-5.png}
		\caption[An examle traindata. RGB, depth and the ground truth]{An examle traindata. RGB, depth and the ground truth}
	\end{figure}
	\fi
	
	\FloatBarrier
	% \clearpage
	One example image for the synthetic dataset with \textbf{80 shapes} and \textbf{160 textures}:
	\begin{figure}[H]
		\centering
		\includegraphics[width=\textwidth]{anhang/example-80-160-1.png}
		\caption[An examle traindata. RGB, depth and the ground truth]{An examle traindata. RGB, depth and the ground truth}
	\end{figure}
	\iffalse
	\begin{figure}[H]
		\centering
		\includegraphics[width=\textwidth]{anhang/example-80-160-2.png}
		\caption[An examle traindata. RGB, depth and the ground truth]{An examle traindata. RGB, depth and the ground truth}
	\end{figure}
	\begin{figure}[H]
		\centering
		\includegraphics[width=\textwidth]{anhang/example-80-160-3.png}
		\caption[An examle traindata. RGB, depth and the ground truth]{An examle traindata. RGB, depth and the ground truth}
	\end{figure}
	\begin{figure}[H]
		\centering
		\includegraphics[width=\textwidth]{anhang/example-80-160-4.png}
		\caption[An examle traindata. RGB, depth and the ground truth]{An examle traindata. RGB, depth and the ground truth}
	\end{figure}
	\begin{figure}[H]
		\centering
		\includegraphics[width=\textwidth]{anhang/example-80-160-5.png}
		\caption[An examle traindata. RGB, depth and the ground truth]{An examle traindata. RGB, depth and the ground truth}
	\end{figure}
	\fi
	
	\FloatBarrier
	% \clearpage
	One example image for the synthetic dataset with \textbf{160 shapes} and \textbf{10 textures}:
	\begin{figure}[H]
		\centering
		\includegraphics[width=\textwidth]{anhang/example-160-10-1.png}
		\caption[An examle traindata. RGB, depth and the ground truth]{An examle traindata. RGB, depth and the ground truth}
	\end{figure}
	\iffalse
	\begin{figure}[H]
		\centering
		\includegraphics[width=\textwidth]{anhang/example-160-10-2.png}
		\caption[An examle traindata. RGB, depth and the ground truth]{An examle traindata. RGB, depth and the ground truth}
	\end{figure}
	\begin{figure}[H]
		\centering
		\includegraphics[width=\textwidth]{anhang/example-160-10-3.png}
		\caption[An examle traindata. RGB, depth and the ground truth]{An examle traindata. RGB, depth and the ground truth}
	\end{figure}
	\begin{figure}[H]
		\centering
		\includegraphics[width=\textwidth]{anhang/example-160-10-4.png}
		\caption[An examle traindata. RGB, depth and the ground truth]{An examle traindata. RGB, depth and the ground truth}
	\end{figure}
	\begin{figure}[H]
		\centering
		\includegraphics[width=\textwidth]{anhang/example-160-10-5.png}
		\caption[An examle traindata. RGB, depth and the ground truth]{An examle traindata. RGB, depth and the ground truth}
	\end{figure}
	\fi
	
	\FloatBarrier
	\clearpage
	One example image for the synthetic dataset with \textbf{160 shapes} and \textbf{80 textures}:
	\begin{figure}[H]
		\centering
		\includegraphics[width=\textwidth]{anhang/example-160-80-1.png}
		\caption[An examle traindata. RGB, depth and the ground truth]{An examle traindata. RGB, depth and the ground truth}
	\end{figure}
	\iffalse
	\begin{figure}[H]
		\centering
		\includegraphics[width=\textwidth]{anhang/example-160-80-2.png}
		\caption[An examle traindata. RGB, depth and the ground truth]{An examle traindata. RGB, depth and the ground truth}
	\end{figure}
	\begin{figure}[H]
		\centering
		\includegraphics[width=\textwidth]{anhang/example-160-80-3.png}
		\caption[An examle traindata. RGB, depth and the ground truth]{An examle traindata. RGB, depth and the ground truth}
	\end{figure}
	\begin{figure}[H]
		\centering
		\includegraphics[width=\textwidth]{anhang/example-160-80-4.png}
		\caption[An examle traindata. RGB, depth and the ground truth]{An examle traindata. RGB, depth and the ground truth}
	\end{figure}
	\begin{figure}[H]
		\centering
		\includegraphics[width=\textwidth]{anhang/example-160-80-5.png}
		\caption[An examle traindata. RGB, depth and the ground truth]{An examle traindata. RGB, depth and the ground truth}
	\end{figure}
	\fi
	
	\FloatBarrier
	% \clearpage
	One example image for the synthetic dataset with \textbf{160 shapes} and \textbf{160 textures}:
	\begin{figure}[H]
		\centering
		\includegraphics[width=\textwidth]{anhang/example-160-160-1.png}
		\caption[An examle traindata. RGB, depth and the ground truth]{An examle traindata. RGB, depth and the ground truth}
	\end{figure}
	\iffalse
	\begin{figure}[H]
		\centering
		\includegraphics[width=\textwidth]{anhang/example-160-160-2.png}
		\caption[An examle traindata. RGB, depth and the ground truth]{An examle traindata. RGB, depth and the ground truth}
	\end{figure}
	\begin{figure}[H]
		\centering
		\includegraphics[width=\textwidth]{anhang/example-160-160-3.png}
		\caption[An examle traindata. RGB, depth and the ground truth]{An examle traindata. RGB, depth and the ground truth}
	\end{figure}
	\begin{figure}[H]
		\centering
		\includegraphics[width=\textwidth]{anhang/example-160-160-4.png}
		\caption[An examle traindata. RGB, depth and the ground truth]{An examle traindata. RGB, depth and the ground truth}
	\end{figure}
	\begin{figure}[H]
		\centering
		\includegraphics[width=\textwidth]{anhang/example-160-160-5.png}
		\caption[An examle traindata. RGB, depth and the ground truth]{An examle traindata. RGB, depth and the ground truth}
	\end{figure}
	\fi


\chapter{Shape vs. Texture Attention Test Data}
\label{appendix:testdata-examples-bias}

	This appendix presents samples of the handcrafted test images for the Shape vs. Texture Attention Test \ref{sec:shape-texutre-attention-test}.\\
	The test data is separated into images with multiple textures and only one texture for all objects. There is an equal distribution of images using known textures and known shapes, known textures and unknown shapes, unknown textures and known shapes, and unknown textures and unknown shapes.\\
	Starting with \textbf{multiple} textures on every object.
	
	\begin{figure}[H]
		\centering
		\includegraphics[width=0.9\textwidth]{anhang/bias/3xM_bias_test_multiple_1.png}
		\caption[An examle testdata for Shape vs. Texture Attention Test with multiple textures. RGB and depth images]{An examle testdata for Shape vs. Texture Attention Test with multiple textures. RGB and depth images}
	\end{figure}
	\begin{figure}[H]
		\centering
		\includegraphics[width=0.9\textwidth]{anhang/bias/3xM_bias_test_multiple_2.png}
		\caption[An examle testdata for Shape vs. Texture Attention Test with multiple textures. RGB and depth images]{An examle testdata for Shape vs. Texture Attention Test with multiple textures. RGB and depth images}
	\end{figure}
	\begin{figure}[H]
		\centering
		\includegraphics[width=0.9\textwidth]{anhang/bias/3xM_bias_test_multiple_3.png}
		\caption[An examle testdata for Shape vs. Texture Attention Test with multiple textures. RGB and depth images]{An examle testdata for Shape vs. Texture Attention Test with multiple textures. RGB and depth images}
	\end{figure}
	\begin{figure}[H]
		\centering
		\includegraphics[width=0.9\textwidth]{anhang/bias/3xM_bias_test_multiple_4.png}
		\caption[An examle testdata for Shape vs. Texture Attention Test with multiple textures. RGB and depth images]{An examle testdata for Shape vs. Texture Attention Test with multiple textures. RGB and depth images}
	\end{figure}
	
	\clearpage
	Test data where all shapes have the same texture.
	
	\begin{figure}[H]
		\centering
		\includegraphics[width=0.9\textwidth]{anhang/bias/3xM_bias_test_one_1.png}
		\caption[An examle testdata for Shape vs. Texture Attention Test with one textures RGB and depth images]{An examle testdata for Shape vs. Texture Attention Test with one textures RGB and depth images}
	\end{figure}
	\begin{figure}[H]
		\centering
		\includegraphics[width=0.9\textwidth]{anhang/bias/3xM_bias_test_one_2.png}
		\caption[An examle testdata for Shape vs. Texture Attention Test with one textures RGB and depth images]{An examle testdata for Shape vs. Texture Attention Test with one textures RGB and depth images}
	\end{figure}
	\begin{figure}[H]
		\centering
		\includegraphics[width=0.9\textwidth]{anhang/bias/3xM_bias_test_one_3.png}
		\caption[An examle testdata for Shape vs. Texture Attention Test with one textures RGB and depth images]{An examle testdata for Shape vs. Texture Attention Test with one textures RGB and depth images}
	\end{figure}
	\begin{figure}[H]
		\centering
		\includegraphics[width=0.9\textwidth]{anhang/bias/3xM_bias_test_one_4.png}
		\caption[An examle testdata for Shape vs. Texture Attention Test with one textures RGB and depth images]{An examle testdata for Shape vs. Texture Attention Test with one textures RGB and depth images}
	\end{figure}
	
	
	
	
	
	
	
	
	
	
	

\chapter{In-Distribution Performance and Generalization Test Data Examples}
\label{appendix:testdata-examples-in-distribution}

	This appendix shows 5 example data points for every of the 4 test datasets from the In-Distribution Performance and Generalization Test \ref{sec:in-distribution-performance-generalization}.
	
	Five examples from in-distribution testdata \ref{sec:in-distribution-performance-generalization} for \textbf{known shapes} and \textbf{known textures}:
	\begin{figure}[H]
		\centering
		\includegraphics[width=\textwidth]{anhang/test-example-known_known-1.png}
		\caption[An examle testdata for in-distribution performance experiment. RGB, depth and mask images]{An examle testdata for in-distribution performance experiment. RGB, depth and mask images}
	\end{figure}
	\begin{figure}[H]
		\centering
		\includegraphics[width=\textwidth]{anhang/test-example-known_known-2.png}
		\caption[An examle testdata for in-distribution performance experiment. RGB, depth and mask images]{An examle testdata for in-distribution performance experiment. RGB, depth and mask images}
	\end{figure}
	\begin{figure}[H]
		\centering
		\includegraphics[width=\textwidth]{anhang/test-example-known_known-3.png}
		\caption[An examle testdata for in-distribution performance experiment. RGB, depth and mask images]{An examle testdata for in-distribution performance experiment. RGB, depth and mask images}
	\end{figure}
	\begin{figure}[H]
		\centering
		\includegraphics[width=\textwidth]{anhang/test-example-known_known-4.png}
		\caption[An examle testdata for in-distribution performance experiment. RGB, depth and mask images]{An examle testdata for in-distribution performance experiment. RGB, depth and mask images}
	\end{figure}
	\begin{figure}[H]
		\centering
		\includegraphics[width=\textwidth]{anhang/test-example-known_known-5.png}
		\caption[An examle testdata for in-distribution performance experiment. RGB, depth and mask images]{An examle testdata for in-distribution performance experiment. RGB, depth and mask images}
	\end{figure}
	
	
	\FloatBarrier
	% \vspace{3cm}
	\clearpage
	Five examples from in-distribution testdata \ref{sec:in-distribution-performance-generalization} for \textbf{known shapes} and \textbf{unknow textures}:
	\begin{figure}[H]
		\centering
		\includegraphics[width=\textwidth]{anhang/test-example-known_unknown-1.png}
		\caption[An examle testdata for in-distribution performance experiment. RGB, depth and mask images]{An examle testdata for in-distribution performance experiment. RGB, depth and mask images}
	\end{figure}
	\begin{figure}[H]
		\centering
		\includegraphics[width=\textwidth]{anhang/test-example-known_unknown-2.png}
		\caption[An examle testdata for in-distribution performance experiment. RGB, depth and mask images]{An examle testdata for in-distribution performance experiment. RGB, depth and mask images}
	\end{figure}
	\begin{figure}[H]
		\centering
		\includegraphics[width=\textwidth]{anhang/test-example-known_unknown-3.png}
		\caption[An examle testdata for in-distribution performance experiment. RGB, depth and mask images]{An examle testdata for in-distribution performance experiment. RGB, depth and mask images}
	\end{figure}
	\begin{figure}[H]
		\centering
		\includegraphics[width=\textwidth]{anhang/test-example-known_unknown-4.png}
		\caption[An examle testdata for in-distribution performance experiment. RGB, depth and mask images]{An examle testdata for in-distribution performance experiment. RGB, depth and mask images}
	\end{figure}
	\begin{figure}[H]
		\centering
		\includegraphics[width=\textwidth]{anhang/test-example-known_unknown-5.png}
		\caption[An examle testdata for in-distribution performance experiment. RGB, depth and mask images]{An examle testdata for in-distribution performance experiment. RGB, depth and mask images}
	\end{figure}
	
	
	
	\FloatBarrier
	% \vspace{3cm}
	\clearpage
	Five examples from in-distribution testdata \ref{sec:in-distribution-performance-generalization} for \textbf{unknown shapes} and \textbf{known textures}:
	\begin{figure}[H]
		\centering
		\includegraphics[width=\textwidth]{anhang/test-example-unknown_known-1.png}
		\caption[An examle testdata for in-distribution performance experiment. RGB, depth and mask images]{An examle testdata for in-distribution performance experiment. RGB, depth and mask images}
	\end{figure}
	\begin{figure}[H]
		\centering
		\includegraphics[width=\textwidth]{anhang/test-example-unknown_known-2.png}
		\caption[An examle testdata for in-distribution performance experiment. RGB, depth and mask images]{An examle testdata for in-distribution performance experiment. RGB, depth and mask images}
	\end{figure}
	\begin{figure}[H]
		\centering
		\includegraphics[width=\textwidth]{anhang/test-example-unknown_known-3.png}
		\caption[An examle testdata for in-distribution performance experiment. RGB, depth and mask images]{An examle testdata for in-distribution performance experiment. RGB, depth and mask images}
	\end{figure}
	\begin{figure}[H]
		\centering
		\includegraphics[width=\textwidth]{anhang/test-example-unknown_known-4.png}
		\caption[An examle testdata for in-distribution performance experiment. RGB, depth and mask images]{An examle testdata for in-distribution performance experiment. RGB, depth and mask images}
	\end{figure}
	\begin{figure}[H]
		\centering
		\includegraphics[width=\textwidth]{anhang/test-example-unknown_known-5.png}
		\caption[An examle testdata for in-distribution performance experiment. RGB, depth and mask images]{An examle testdata for in-distribution performance experiment. RGB, depth and mask images}
	\end{figure}
	
	
	
	\FloatBarrier
	% \vspace{3cm}
	\clearpage
	Five examples from in-distribution testdata \ref{sec:in-distribution-performance-generalization} for \textbf{unknown shapes} and \textbf{unknown textures}:
	\begin{figure}[H]
		\centering
		\includegraphics[width=\textwidth]{anhang/test-example-unknown_unknown-1.png}
		\caption[An examle testdata for in-distribution performance experiment. RGB, depth and mask images]{An examle testdata for in-distribution performance experiment. RGB, depth and mask images}
	\end{figure}
	\begin{figure}[H]
		\centering
		\includegraphics[width=\textwidth]{anhang/test-example-unknown_unknown-2.png}
		\caption[An examle testdata for in-distribution performance experiment. RGB, depth and mask images]{An examle testdata for in-distribution performance experiment. RGB, depth and mask images}
	\end{figure}
	\begin{figure}[H]
		\centering
		\includegraphics[width=\textwidth]{anhang/test-example-unknown_unknown-3.png}
		\caption[An examle testdata for in-distribution performance experiment. RGB, depth and mask images]{An examle testdata for in-distribution performance experiment. RGB, depth and mask images}
	\end{figure}
	\begin{figure}[H]
		\centering
		\includegraphics[width=\textwidth]{anhang/test-example-unknown_unknown-4.png}
		\caption[An examle testdata for in-distribution performance experiment. RGB, depth and mask images]{An examle testdata for in-distribution performance experiment. RGB, depth and mask images}
	\end{figure}
	\begin{figure}[H]
		\centering
		\includegraphics[width=\textwidth]{anhang/test-example-unknown_unknown-5.png}
		\caption[An examle testdata for in-distribution performance experiment. RGB, depth and mask images]{An examle testdata for in-distribution performance experiment. RGB, depth and mask images}
	\end{figure}
	

	

\chapter{Optonic Bin-Picking Dataset Examples}
\label{appendix:testdata-examples-simtoreal}

	This study proposes a real-world dataset for instance-segmentation as described in section \ref{sec:test-data}. Here are samples from this dataset with RGB, depth and (colored) mask.
	
	\begin{figure}[H]
		\centering
		\includegraphics[width=\textwidth]{anhang/real/real_data_sample_1.png}
		\caption[An sample observation from the sim-to-real test dataset.]{An sample observation from the sim-to-real test dataset}
	\end{figure}
	\begin{figure}[H]
		\centering
		\includegraphics[width=\textwidth]{anhang/real/real_data_sample_2.png}
		\caption[An sample observation from the sim-to-real test dataset.]{An sample observation from the sim-to-real test dataset}
	\end{figure}
	\begin{figure}[H]
		\centering
		\includegraphics[width=\textwidth]{anhang/real/real_data_sample_3.png}
		\caption[An sample observation from the sim-to-real test dataset.]{An sample observation from the sim-to-real test dataset}
	\end{figure}
	\begin{figure}[H]
		\centering
		\includegraphics[width=\textwidth]{anhang/real/real_data_sample_4.png}
		\caption[An sample observation from the sim-to-real test dataset.]{An sample observation from the sim-to-real test dataset}
	\end{figure}
	\begin{figure}[H]
		\centering
		\includegraphics[width=\textwidth]{anhang/real/real_data_sample_5.png}
		\caption[An sample observation from the sim-to-real test dataset.]{An sample observation from the sim-to-real test dataset}
	\end{figure}
	\begin{figure}[H]
		\centering
		\includegraphics[width=\textwidth]{anhang/real/real_data_sample_6.png}
		\caption[An sample observation from the sim-to-real test dataset.]{An sample observation from the sim-to-real test dataset}
	\end{figure}
	\begin{figure}[H]
		\centering
		\includegraphics[width=\textwidth]{anhang/real/real_data_sample_7.png}
		\caption[An sample observation from the sim-to-real test dataset.]{An sample observation from the sim-to-real test dataset}
	\end{figure}
	\begin{figure}[H]
		\centering
		\includegraphics[width=\textwidth]{anhang/real/real_data_sample_8.png}
		\caption[An sample observation from the sim-to-real test dataset.]{An sample observation from the sim-to-real test dataset}
	\end{figure}
	\begin{figure}[H]
		\centering
		\includegraphics[width=\textwidth]{anhang/real/real_data_sample_9.png}
		\caption[An sample observation from the sim-to-real test dataset.]{An sample observation from the sim-to-real test dataset}
	\end{figure}
	\begin{figure}[H]
		\centering
		\includegraphics[width=\textwidth]{anhang/real/real_data_sample_10.png}
		\caption[An sample observation from the sim-to-real test dataset.]{An sample observation from the sim-to-real test dataset}
	\end{figure}
	\begin{figure}[H]
		\centering
		\includegraphics[width=\textwidth]{anhang/real/real_data_sample_11.png}
		\caption[An sample observation from the sim-to-real test dataset.]{An sample observation from the sim-to-real test dataset}
	\end{figure}
	\begin{figure}[H]
		\centering
		\includegraphics[width=\textwidth]{anhang/real/real_data_sample_12.png}
		\caption[An sample observation from the sim-to-real test dataset.]{An sample observation from the sim-to-real test dataset}
	\end{figure}
	\begin{figure}[H]
		\centering
		\includegraphics[width=\textwidth]{anhang/real/real_data_sample_13.png}
		\caption[An sample observation from the sim-to-real test dataset.]{An sample observation from the sim-to-real test dataset}
	\end{figure}
	\begin{figure}[H]
		\centering
		\includegraphics[width=\textwidth]{anhang/real/real_data_sample_14.png}
		\caption[An sample observation from the sim-to-real test dataset.]{An sample observation from the sim-to-real test dataset}
	\end{figure}
	\begin{figure}[H]
		\centering
		\includegraphics[width=\textwidth]{anhang/real/real_data_sample_15.png}
		\caption[An sample observation from the sim-to-real test dataset.]{An sample observation from the sim-to-real test dataset}
	\end{figure}

	


\end{document}
